
rated by a family \((p_{\iota})\) of elements such that each \(A p_{\iota}\) is a prime ideal (note that every divisorial prime ideal which meets S is principal).

¶ 15. (a) Let A be an integral domain and a, b two elements ≠0 of A such that \(Aa \cap Ab = Aab\) . Show that in the polynomial ring A{[}X{]} the principal ideal generated by \(aX + b\) is prime (prove that every polynomial \(f(X) \in A[X]\) such that \(f(-b/a) = 0\) is a multiple of \(aX + b\) , arguing induction on the degree of f).

\begin{enumerate}
\def\labelenumi{(\alph{enumi})}
\setcounter{enumi}{1}
\item
  Let A be a Krull domain and \(a, b\) two elements of A such that \(Aa\) and \(Aa + Ab\) are prime and distinct. Show that \(\mathrm{A}[\mathrm{X}] / (a\mathrm{X} + b)\) is a Krull domain and that \(\mathrm{C}(\mathrm{A}[\mathrm{X}] / (a\mathrm{X} + b))\) is isomorphic to \(\mathrm{C}(\mathrm{A})\) . (Note, with the aid of (a), that \(\mathrm{A}[\mathrm{X}] / (a\mathrm{X} + b)\) is isomorphic to \(\mathrm{A}[b / a] = \mathrm{B}\) ; show that \(Ba\) is prime in B. Observe that \(\mathrm{A}[a^{-1}] = \mathrm{B}[a^{-1}]\) is a Krull domain and use Exercise 12.)
\item
  Let A be an integrally closed Noetherian domain and a, b two elements of the Jacobson radical r of A. Suppose that the ideal \(Aa + Ab\) is a non-divisorial prime ideal; show that Aa and Ab are prime ideals. (Prove first that \(Aa \cap Ab = Aab\) considering the ideals of \(\operatorname{Ass}(A/Aa)\) ). Show secondly that the relations \(xy \in Ab\) and \(y \notin Aa + Ab\) imply \(x \in Ab + Aa^{h}\) for all h, by induction on h; deduce that then \(x \in Ab\) . Finally, deduce that, if \(xy \in Ab\) , \(x \notin Ab\) , \(y \notin Ab\) , then necessarily \(x \in Ab + Aa^{h}\) and \(y \in Ab + Aa^{h}\) for all h, arguing by induction on h; deduce a contradiction.)
\end{enumerate}

\begin{enumerate}
\def\labelenumi{\arabic{enumi}.}
\setcounter{enumi}{15}
\tightlist
\item
  Let \(A = \bigoplus_{n \in Z} A_n\) be a graded Krull domain. Let \(\mathrm{D}_g(\mathrm{A})\) (resp. \(\mathrm{F}_g(\mathrm{A})\) ) denote the group generated by the divisors of the graded integral divisorial ideals (resp. by the divisors \(\operatorname{div}(a)\) , where a is homogeneous in A); write \(\mathrm{C}_g(\mathrm{A}) = \mathrm{D}_g(\mathrm{A}) / \mathrm{F}_g(\mathrm{A})\) . Let S be the multiplicative subset of homogeneous elements \(\neq 0\) of A.
\end{enumerate}

\begin{enumerate}
\def\labelenumi{(\alph{enumi})}
\item
  Show that every prime ideal of A of height 1 which meets S is graded (use Chapter III, § 1, no. 4, Proposition 5).
\item
  Let \(\mathbf{B} = \mathbf{S}^{-1}\mathbf{A}\) , which is a graded ring (Chapter II, § 2, no. 9); write \(\mathbf{B} = \bigoplus_{n\in \mathbf{Z}}\mathbf{B}_n\) ; show that \(\mathbf{B}_0\) is a field and that, if \(\mathbf{A}\neq \mathbf{A}_0\) , \(\mathbf{B}\) is isomorphic to the ring \(\mathbf{B}_0[\mathbf{X},\mathbf{X}^{-1}]\) of rational functions \(\mathrm{P}(\mathbf{X}) / \mathbf{X}^k\) , where \(\mathrm{P}(\mathbf{X})\in \mathrm{B}_0[\mathbf{X}]\) .
\item
  Show that \(C_{g}(A)\) and \(C(A)\) are canonically isomorphic (use (a), (b) and equation (5) of no. 10).
\end{enumerate}

\begin{enumerate}
\def\labelenumi{\arabic{enumi}.}
\setcounter{enumi}{16}
\tightlist
\item
  Let \(\mathbf{A} = \bigoplus_{n\in \mathbf{Z}}\mathbf{A}_n\) be a graded Krull domain and \(p\in \mathbf{A}_1\) be such that \(\mathbf{A}p\) is prime and \(\neq 0\) .
\end{enumerate}

\begin{enumerate}
\def\labelenumi{(\alph{enumi})}
\item
  Show that, if B is the ring \(\mathbf{A}_{(p)}\) defined in Chapter III, § 1, Exercise 1 (a), B is a Krull domain and the groups C(A) and C(B) are canonically isomorphic. (Show first that C(B) is isomorphic to C(B{[}p,p \(^{-1}\) {]}) and observe that B{[}p,p \(^{-1}\) {]} = A{[}p \(^{-1}\) {]}.)
\item
  Under the hypotheses of Exercise 15 (b), show that the groups \(\mathbf{C}(\mathbf{A})\) and \(\mathbf{C}(\mathbf{A}[\mathbf{X},\mathbf{Y}] / (a\mathbf{X} + b\mathbf{Y}))\) are isomorphic.
\end{enumerate}

¶ 18. Let \(A = \bigoplus_{n \in \mathbf{Z}} A_n\) be a graded Krull domain with positive degrees, such that \(A_0\) is a field and let \(m = \bigoplus_{n \geqslant 1} A_n\) , a maximal ideal of \(A\) . Let \(S\) be the multiplicative set of homogeneous elements \(\neq 0\) of \(A\) , so that the ring \(S^{-1}A\) is a principal ideal domain (Exercise 16 (b)).

\begin{enumerate}
\def\labelenumi{(\alph{enumi})}
\item
  Let p be a prime ideal of A of height 1 which meets A --- m; then p is not graded (otherwise p = A) and the ideal S \(^{-1}\) p of S \(^{-1}\) A is principal, generated by an element of the form x = 1 + x \(_{1}\) + · · · + x \(_{n}\) where x \(_{j}\) ∈ (S \(^{-1}\) A) \(_{j}\) . Writing an element of p of the form 1 + a \(_{1}\) + · · · + a \(_{q}\) (a \(_{i}\) ∈ A \(_{i}\) ) as a multiple of x in S \(^{-1}\) A and using Chapter V, § 1, no. 3, Proposition 11, show that x ∈ p; writing finally every element of p as a multiple of x in S \(^{-1}\) A, show that p = Ax.
\item
  Deduce from (a) that C(A) and C(A \(_{m}\) ) are canonically isomorphic (use Proposition 17 of no. 10).
\end{enumerate}

\begin{enumerate}
\def\labelenumi{\arabic{enumi}.}
\setcounter{enumi}{18}
\tightlist
\item
  Let A be a local Krull domain and m its maximal ideal; if \(\mathbf{A}' = (\mathbf{A}[\mathbf{X}])_{\mathfrak{m}\mathbf{A}[\mathbf{X}]}\) , show that \(\mathbf{C}(\mathbf{A})\) and \(\mathbf{C}(\mathbf{A}')\) are canonically isomorphic. (Apply the criterion of Proposition 17 of no. 10, using Exercise 12 (c).)
\end{enumerate}

¶ 20. An integral domain A is called Bezoutian (or a Bezout domain) if every finitely generated ideal of A is principal. Every Bezoutian Noetherian domain is a principal ideal domain. Every valuation ring is a Bezoutian domain (and therefore a Bezoutian domain is not therefore necessarily Noetherian nor completely integrally closed).

\begin{enumerate}
\def\labelenumi{(\alph{enumi})}
\item
  Show that every Bezoutian domain is integrally closed (cf.~Exercise 6). If an integral domain is the union of a right directed family of Bezoutian subdomains, it is Bezoutian. If A is a Bezoutian domain, so is S \(^{-1}\) A for every multiplicative subset S of A such that 0 \(\notin\) S.
\item
  Let \(v_{i}\) ( \(1 \leqslant i \leqslant n\) ) be independent valuations on a field \(K\) and \(A_{i}\) the ring of the valuation \(v_{i}\) . Show that the intersection \(A\) of the \(A_{i}\) is a Bezoutian domain. (If an ideal \(a\) of \(A\) is finitely generated, the set of \(v_{i}(x)\) for \(x \in a\) admits a least element \(\alpha_{i}\) in the value group of \(v_{i}\) ; for all \(i\) , let \(x_{i} \in a\) be such that \(v_{i}(x_{i}) = \alpha_{i}\) . Using the Approximation Theorem (Chapter VI, §7, no. 2,
\end{enumerate}

Theorem 1), show that there are elements \(a_i \in \mathbf{A}\) such that \(x = \sum_{i=1}^{n} a_i x_i \in \mathfrak{a}\) satisfies the relations \(v_i(x) = \alpha_i\) for \(1 \leqslant i \leqslant n\) . If the \(v_i\) are discrete valuations, \(\mathbf{A}\) is a principal ideal domain.

\begin{enumerate}
\def\labelenumi{(\alph{enumi})}
\setcounter{enumi}{2}
\tightlist
\item
  Let K be an algebraically closed field of characteristic \(\neq2\) and let A be the subring of an algebraic closure of \(\mathbf{K}(\mathbf{X})\) generated by \(\mathbf{K}\) and two sequences of elements \((x_{n})\) , \((1 / x_{n})\) ( \(1 \leqslant n < +\infty\) ), where \(x_{1} = \mathbf{X}\) and \(x_{n-1} = x_{n}^{2}\) . Show that \(\mathbf{A}\) is Bezoutian (use (a)). If \(\mathfrak{p}\) is a prime ideal \(\neq 0\) of \(\mathbf{A}\) , show that \(\mathfrak{p}\) is generated by a sequence of elements of the form \(x_{n} - a_{n}\) , where \(a_{n} \in \mathbf{K}\) , \(a_{n} \neq 0\) and \(a_{n}^{2} = a_{n-1}\) (consider for all \(n\) the intersection \(\mathfrak{p} \cap \mathbf{K}[x_{n}, 1 / x_{n}]\) ). Show that \(\mathfrak{p}\) is not finitely generated.
\end{enumerate}

\begin{enumerate}
\def\labelenumi{\arabic{enumi}.}
\setcounter{enumi}{20}
\tightlist
\item
  An integral domain A is called pseudo-Bezoutian (resp. pseudo-principal) if, in the notation of Exercise 5, the group K*/U is lattice-ordered (completely lattice ordered). Every Bezoutian * (resp. factorial) * domain is pseudo-Bezoutian * (resp. pseudo-principal) *. Every pseudo-principal domain is pseudo-Bezoutian. A valuation ring whose order group is R is pseudo-principal but not a principal ideal domain. Every pseudo-Bezoutian (resp. pseudo-principal) domain is integrally closed (resp. completely integrally closed) (use Exercise 6). * Every Noetherian pseudo-Bezoutian domain is factorial. * Give an example of an integrally closed Noetherian (and therefore Krull) domain which is not pseudo-Bezoutian (cf.~Exercise 2). If A is pseudo-Bezoutian, so is S \(^{-1}\) A for every multiplicative subset S of A such that 0∉S (use Chapter II, § 2, Exercise 1).
\end{enumerate}

¶ 22. (a) Let \(\Gamma\) be a lattice-ordered additive group and A the algebra of \(\Gamma\) over a field \(k\) ; A is an integral domain (Algebra, Chapter II, §11, no. 4, Proposition 8). Every element \(x \neq 0\) in A may be written uniquely as \(x = \sum_{i=1}^{n} \alpha_i e^{\nu_i}\) , where the \(\nu_i\) are distinct elements of \(\Gamma\) , the \(e^{\nu_i}\) the corresponding elements of the canonical basis of A over \(k\) and the \(\alpha_i\) elements \(\neq 0\) of \(k\) . Write \(\phi(x) = \inf(\nu_1, \ldots, \nu_n)\) in \(\Gamma\) ; show that \(\phi(x + y) \geqslant \inf(\phi(x), \phi(y))\) if \(x, y\) and \(x + y\) are \(\neq 0\) in A and \(\phi(xy) = \phi(x) + \phi(y)\) if \(xy \neq 0\) in A. (To show the second assertion, establish first the following lemma: given a finite family \((\xi_j)\) of elements of \(\Gamma\) , in order that \(\inf_{i} (\xi_j) = 0\) , it is necessary and sufficient that for all \(\eta > 0\) in \(\Gamma\) there exist an index \(j\) and an element \(\zeta \in \Gamma\) such that \(0 < \zeta \leqslant \eta\) and \(\inf(\xi_j, \zeta) = 0\) . Apply this lemma reducing it to the case where \(\phi(x) = \phi(y) = 0\) .)

\begin{enumerate}
\def\labelenumi{(\alph{enumi})}
\setcounter{enumi}{1}
\tightlist
\item
  Deduce from (a) that, if \(\mathbf{K}\) is the field of fractions of \(\mathbf{A},\phi\) can be extended to a homomorphism from \(\mathbf{K}^*\) to \(\Gamma\) (also denoted by \(\phi\) ) such that
\end{enumerate}

\[
\phi (x + y) \geqslant \inf (\phi (x), \phi (y))
\]

if \(x, y\) and \(x + y\) are \(\neq 0\) in \(\mathbf{K}\) . Deduce that, if \(\mathbf{B}\) is the set of \(x \in \mathbf{K}\) such that \(x = 0\) or \(\phi(x) \geqslant 0\) , \(\mathbf{B}\) is a domain whose field of fractions is \(\mathbf{K}\) and such that, if \(\mathbf{U}\) is the group of invertible elements of \(\mathbf{B}\) , \(\mathbf{K}^*/\mathbf{U}\) is isomorphic to \(\Gamma\) (and in particular \(\mathbf{B}\) is a pseudo-Bezoutian domain). Derive examples of a pseudo-Bezoutian domain which is not completely integrally closed, a completely integrally closed pseudo-Bezoutian domain which is not pseudo-principal (cf.

Algebra, Chapter VI, § 1, Exercise 31) * and a pseudo-principal domain which is not factorial. *

* (c) Derive from (b) an example of an integral domain which is the union of a right directed family of factorial subdomains and is completely integrally closed but not pseudo-Bezoutian. (Let \(\theta\) be an irrational number in {[}0, 1{]} and for every integer \(j\) let \(q_j\) be the greatest integer such that \(q_j / 2^j < \theta\) . For all \(j\) define on the product group \(\mathbf{Q} \times \mathbf{Q}\) a lattice-ordered group structure by taking the set \((\mathrm{G}_j)_+\) of elements \(\geqslant 0\) of this group to be the ordered pairs \((\xi, \eta)\) such that \(\xi \geqslant 0\) and \(0 \leqslant \eta \leqslant (q_j 2^{-j}) \xi\) . *

\begin{enumerate}
\def\labelenumi{\arabic{enumi}.}
\setcounter{enumi}{22}
\tightlist
\item
  \begin{enumerate}
  \def\labelenumii{(\alph{enumii})}
  \tightlist
  \item
    Let A be a pseudo-Bezoutian domain (Exercise 21) and K its field of fractions; for every polynomial \(f \in \mathrm{A}[\mathrm{X}]\) a g.c.d. of the coefficients of \(f\) (determined to within an invertible element of A) is called a content of \(f\) . Let \(f, g\) be two polynomials of \(\mathrm{A}[\mathrm{X}]\) ; show that for \(f\) to divide \(g\) in \(\mathrm{A}[\mathrm{X}]\) , it is necessary and sufficient that \(f\) divide \(g\) in \(\mathrm{K}[\mathrm{X}]\) and that a content of \(f\) divide a content of \(g\) (use Exercise 12) (cf.~Exercise 30 (c)).
  \end{enumerate}
\end{enumerate}

\begin{enumerate}
\def\labelenumi{(\alph{enumi})}
\setcounter{enumi}{1}
\item
  Deduce from (a) that, if A is a pseudo-Bezoutian (resp. pseudo-principal) domain, then so is A{[}X{]}. Moreover, if \((a_{t})\) is a finite (resp. any) family of elements \(\neq0\) of A and d a g.c.d. of this family in A, d is also a g.c.d. of this family in A{[}X{]}.
\item
  Deduce from (b) that, if A is pseudo-Bezoutian (resp. pseudo-principal) then \(A[X_{\lambda}]_{\lambda \in L}\) is pseudo-Bezoutian (resp. pseudo-principal) for every family \((\mathbf{X}_{\lambda})_{\lambda \in L}\) of indeterminates.
\end{enumerate}

\begin{enumerate}
\def\labelenumi{\arabic{enumi}.}
\setcounter{enumi}{23}
\item
  Let K be an algebraically closed field and A = K{[}X, Y{]} the polynomial ring in two indeterminates over K, which is a Krull domain * (and even a factorial domain). * For every ordered pair \((\alpha, \beta) \in \mathrm{K}^{2}\) let \(w_{\alpha,\beta}\) be the discrete valuation on A such that, for every polynomial \(f \neq 0\) , \(w_{\alpha,\beta}(f)\) is the least degree of the monomials \(\neq 0\) in \(f(\mathrm{X} + \alpha, \mathrm{Y} + \beta)\) . Show that A is the intersection of the rings of the valuations \(w_{\alpha,\beta}\) , none of which is an essential valuation of A.
\item
  An integrally closed domain \(\mathbf{A}\) is said to be of finite character if there exists a family \((v_{\iota})_{\iota \in I}\) of valuations on the field of fractions of \(\mathbf{A}\) satisfying properties \((\mathrm{AK}_{\Pi})\) and \((\mathrm{AK}_{\mathrm{III}})\) .
\end{enumerate}

\begin{enumerate}
\def\labelenumi{(\alph{enumi})}
\tightlist
\item
  Show that, if A is integrally closed and of finite character, for every element \(x \neq 0\) in A, there can only be a finite number of extremal elements of the ordered group of principal divisors which are \(\leqslant \operatorname{div}(x)\) . (Let \(J \subset I\) be the finite set of indices such that \(v_{\iota}(x) > 0\) . If \(m\) is the number of elements in J, show that there cannot be \(m + 1\) elements \(y_{j}\) ( \(1 \leqslant j \leqslant m + 1\) ) dividing \(x\) and such that the elements \(\operatorname{div}(y_{j})\) are extremal, observing that, for all \(j\) , \(\operatorname{div}(y_{j})\) must be prime to \(\sum_{k \neq j} \operatorname{div}(y_{k})\) .)
\end{enumerate}

* (b) Show that the ring of integral functions of one complex variable (Chapter V, § 1, Exercise 12) is a pseudo-principal domain (Exercise 21) but not a domain of finite character (use (a)). *

¶ 26. Let A be an integral domain and K its field of fractions. A valuation v on K is called essential for A if the ring of v is a local ring \(A_{p}\) at a prime ideal p of A (the intersection of A and the ideal of v).

\begin{enumerate}
\def\labelenumi{(\alph{enumi})}
\item
  Let \(v\) be an essential valuation for \(A\) of height 1 and let \(p\) be the prime ideal \(A\) consisting of those \(x\) such that \(v(x) > 0\) . Show that \(p\) is of height 1. If \(x \in K\) does not belong to \(A_p\) , then \(A \cap x^{-1}A \subset p\) . If \(w\) is a valuation on \(K\) whose ring \(B\) contains \(A\) and whose ideal \(q\) satisfies \(q \cap A \subset p\) , show that \(w\) is equivalent to \(v\) .
\item
  An integrally closed domain A is said to be of finite character and real type if there exists a family \((v_{\iota})_{\iota\in I}\) of valuations on K of height 1 satisfying properties \((\mathbf{AK}_{\mathrm{II}})\) and \((\mathbf{AK}_{\mathrm{III}})\) . Show that under these conditions, if \(z\in K\) does not belong to A and p is a prime ideal of A such that \(A\cap z^{-1}A\subset p\) , there exists \(\iota\in I\) such that \(v_{\iota}(z)<0\) and that the prime ideal \(q_{\iota}\) of \(x\in A\) such that \(v_{\iota}(x)>0\) is contained in p.~(Argue by reductio ad absurdum by considering the finite number of \(v_{\iota_{k}}\) such that \(v_{\iota_{k}}(z)<0\) ; prove the existence of an \(a\in A\) such that \(a\notin p\) and \(v_{\iota_{k}}(a)>0\) for all k; deduce that \(a^{n}z\notin A\) for every integer n\textgreater0 and show that this implies a contradiction.) Conclude from this that every essential valuation for A of height 1 is equivalent to one of the \(v_{\iota}\) (use (a)). Every finite intersection of subdomains of K which are of finite character and real type is also a domain of finite character and real type.
\item
  Suppose that the hypotheses of (b) are satisfied and moreover that all the \(v_{\iota}\) are essential. Show that, for every prime ideal p of A of height 1, \(A_{p}\) is the ring of one of the valuations \(v_{\iota}\) (use (b) taking \(z^{-1} \in p\) ). For all \(x \in A\) , the principal ideal Ax then admits a unique reduced primary decomposition (Chapter IV, §2, Exercise 20), the prime ideals corresponding to this decomposition being the prime ideals of height 1 containing x.
\item
  Suppose that the hypotheses of (b) are satisfied. Let S be a multiplicative subset of A not containing 0 and \(J \subset I\) the set of \(\iota \in I\) such that \(v_{\iota}(x) = 0\) on S. Show that the family \((v_{\iota})_{\iota \in I-J}\) satisfies properties \((\mathrm{AK}_{\mathrm{II}})\) and \((\mathrm{AK}_{\mathrm{III}})\) for the ring \(S^{-1}A\) ; this family consists of essential valuations for \(S^{-1}A\) if \((v_{\iota})_{\iota \in I}\) consists of essential valuations for A.
\item
  Generalize Proposition 9 of no. 5 to the case where the hypotheses of (c) are satisfied.
\item
  Suppose the hypotheses of (b) are satisfied. Let \(K'\) be a finite extension of K and \(A'\) the integral closure of A in \(K'\) . Show that the valuations (no two of which are equivalent) on \(K'\) which extend the \(v_{\iota}\) satisfy properties ( \(AK_{II}\) ) and ( \(AK_{III}\) ) for \(A'\) and the latter is therefore a domain of finite character and real type; if the \(v_{\iota}\) are all essential for A, so are their extensions for \(A'\) (argue as in no. 8, Proposition 12 and use also Chapter VI, § 8, no. 3, Remark).
\item
  If \(\mathbf{A}\) is a domain of finite character and real type, so is \(\mathbf{A}[\mathbf{X}]\) ; if the conditions of (c) are fulfilled, determine the essential valuations for A{[}X{]} (argue as in no. 9). Generalize Exercise 8 similarly.
\end{enumerate}

\begin{enumerate}
\def\labelenumi{\arabic{enumi}.}
\setcounter{enumi}{26}
\tightlist
\item
  In Exercise 22 take \(\Gamma\) to be a direct sum of a family \((\Gamma_{\iota})_{\iota\in I}\) , where the \(\Gamma_{\iota}\) are subgroups of the additive group \(\mathbf{R}\) . Show that the ring \(\mathbf{B}\) defined in Exercise 22 (b) is a domain of finite character and real type and that it is the intersection of a family of essential valuation rings for \(\mathbf{B}\) ; moreover, every prime ideal \(\neq 0\) of \(\mathbf{B}\) is maximal.
\end{enumerate}

¶ 28. An integrally closed domain A is said to be of finite character and rational type if there exists a family \((v_{\iota})_{\iota \in I}\) of valuations on its field of fractions K satisfying properties \((\mathbf{AK}_{\mathrm{II}})\) and \((\mathbf{AK}_{\mathrm{III}})\) and whose value groups are subgroups of the additive group Q. Show that the family of essential valuations for A of height 1 satisfy \((\mathbf{AK}_{\mathrm{II}})\) and \((\mathbf{AK}_{\mathrm{III}})\) . (For all \(\iota \in I\) , let \(p_{\iota}\) be the prime ideal of A consisting of the x such that \(v_{\iota}(x) > 0\) and let \(V_{\iota}\) be the ring of the valuation \(v_{\iota}\) . Show that, if there are two indices \(\alpha, \beta\) such that \(p_{\beta} \subset p_{\alpha}\) , then A is the intersection of the \(V_{\iota}\) of index \(\iota \neq \alpha\) . For this, argue by reductio ad absurdum by showing that otherwise there would be an element \(x \in K^{*}\) , an element \(y \in p_{\beta}\) and two positive integers r, s such that \(v_{\beta}(x^{r}y^{s}) > 0\) , \(v_{\alpha}(x^{r}y^{s}) = 0\) and \(v_{\iota}(x^{r}y^{s}) \geqslant 0\) for all \(\iota \in I\) , contrary to the hypothesis.)

\begin{enumerate}
\def\labelenumi{\arabic{enumi}.}
\setcounter{enumi}{28}
\tightlist
\item
  Let A be an integrally closed domain and K its field of fractions.
\end{enumerate}

\begin{enumerate}
\def\labelenumi{(\alph{enumi})}
\item
  Let L be an algebraic extension of K and B the integral closure of A in L. Show that for every fractional ideal a of A, if we write b = aB, \(\tilde{b} \cap A = \tilde{a}\) . (Reduce it to the case where \(A \subset \tilde{a}\) , in other words A: \(a \subset A\) and prove that B: \(b \subset B\) ; for this, let \(x \in B\) : b and let \(c_i\) ( \(1 \leqslant i \leqslant n\) ) be the coefficients of its minimal polynomial over K; note that for all \(y \in a\) the elements \(c_i y^i\) ( \(1 \leqslant i \leqslant n\) ) belong to A and deduce that the \(c_i\) belong to A.)
\item
  Let C be a polynomial ring \(\mathbf{A}[\mathbf{X}_{\lambda}]_{\lambda \in L}\) in any family of indeterminates. Show that for every fractional ideal \(\mathfrak{a}\) of A, if we write \(\mathfrak{c} = \mathfrak{aC}\) , \(\tilde{\mathfrak{c}} \cap \mathbf{A} = \tilde{\mathfrak{a}}\) (same method).
\end{enumerate}

¶ 30. An integral domain is called regularly integrally closed if, in the monoid D(A), every finitely generated divisor (Exercise 11) is a regular element. Every completely integrally closed domain is regularly integrally closed; every pseudo-Bezoutial domain (Exercise 21) is regularly integrally closed.

\begin{enumerate}
\def\labelenumi{(\alph{enumi})}
\item
  If A is regularly integrally closed and \(d, d', d''\) are three finitely generated divisors, the relation \(d + d'' \leqslant d' + d''\) implies \(d \leqslant d'\) .
\item
  Deduce from (a) that, for an integral domain A to be regularly integrally closed, it is necessary and sufficient that, for every fractional ideal a of A such that \(\text{div}(a)\) is a finitely generated divisor, \(a: a = A\) ; in particular A is integrally closed (Exercise 6 (b)). (Use the fact that A: (bc) = (A: b): c for two fractional ideals b, c of A, taking b = a, c = A: a.) Moreover, \(\text{div}(a)\) is then invertible in D(A).
\item
  Let A be a regularly integrally closed domain and K its field of fractions; the monoid \(\mathrm{D}_{f}(\mathrm{A})\) of finitely generated divisors of A generates in \(\mathrm{D}(\mathrm{A})\) a lattice-ordered group \(\mathrm{G}_{f}(\mathrm{A})\) . For every polynomial \(p \in K[X]\) , let \(d(p)\) denote the divisor of the fractional ideal of A generated by the coefficients of p; for every rational function r = p/q of \(\mathrm{K}(\mathrm{X})\) , where p, q are in \(K[X]\) and \(q \neq 0\) , \(d(p) - d(q)\) is well defined on \(\mathrm{G}_{f}(\mathrm{A})\) (independently of the expression of r as a quotient of two polynomials); it is denoted by \(d(r)\) (cf.~Exercise 12 (c)). If we then write \(\gamma(r) = ((r), d(r))\) , where (r) is the fractional principal ideal of \(K[X]\) generated by r, \(\gamma\) is an isomorphism of the ordered group \(\mathcal{P}^{*}(\mathrm{A}[\mathrm{X}])\) (notation of §3, no. 2) onto a subgroup of the product ordered group
\end{enumerate}

\[
\mathscr {P} ^ {*} (\mathrm{K} [ \mathrm{X} ]) \times \mathrm{G} _ {f} (\mathrm{A}) \subset \mathscr {P} ^ {*} (\mathrm{K} [ \mathrm{X} ]) \times \mathrm{D} (\mathrm{A}).
\]

Deduce that \(\mathbf{A}[\mathbf{X}]\) is a regularly integrally closed domain and that \(\mathrm{D}_f(\mathbf{A}[\mathbf{X}])\) is isomorphic to \(\mathcal{P}^* (\mathbf{K}[\mathbf{X}])\times \mathrm{D}_f(\mathbf{A})\) (use (a) and (b) and Exercise 29 (b)).

\begin{enumerate}
\def\labelenumi{(\alph{enumi})}
\setcounter{enumi}{3}
\tightlist
\item
  Let B be an integrally closed domain, K its field of fractions and A the subring of the polynomial ring K{[}X, Y{]} consisting of the polynomials whose constant term belongs to B. Show that, if \(B \neq K\) , A is integrally closed but not regularly integrally closed (show that the divisor \(\inf(\text{div}(X), \text{div}(Y))\) corresponds to a divisorial ideal a such that \(a: a \neq A\) ).
\end{enumerate}

¶ 31. (a) Let A be a regularly integrally closed domain (Exercise 30) and K its field of fractions; for every polynomial P ∈ K{[}X{]}, let c(P) denote the divisor div(a), where a is the fractional ideal of K generated by the coefficients of P. Let B denote the subring of K(X) consisting of the rational functions P/Q such that c(P) ≥ c(Q) (show that this condition depends only on the element P/Q, using Exercise 30 (a) and 12 (c)). Then B ∩ K = A.

\begin{enumerate}
\def\labelenumi{(\alph{enumi})}
\setcounter{enumi}{1}
\item
  Show that B is a Bezoutian domain (Exercise 20) (observe that, if P, Q are two polynomials of A{[}X{]}, P + X \(^{m}\) Q divides P and Q in B if m is sufficiently large).
\item
  Show that, if A is a Krull domain, B is a principal ideal domain (consider an increasing sequence of finitely generated fractional principal ideals of B). Give an example where A is completely integrally closed but B is not a principal ideal domain (cf.~Exercise 22).
\item
  Deduce from (c) an example of a Noetherian domain B for which there exists a subfield K of its field of fractions such that \(B \cap K\) is not Noetherian.
\end{enumerate}

\begin{enumerate}
\def\labelenumi{\arabic{enumi}.}
\setcounter{enumi}{31}
\tightlist
\item
  Let A be an integral domain \((a_{i})_{1 \leqslant i \leqslant n}\) a finite family of elements of A, a the ideal \(\sum_{i} A a_{i}\) and R the submodule of \(A^{n}\) generated by the element \((a_{1}, \ldots, a_{n})\) . Show that, for the torsion submodule of the A-module \(M = A^{n}/R\) to be a direct factor of M, it is necessary and sufficient that
\end{enumerate}

\[
\mathfrak {a} (\mathrm{A}: \mathfrak {a}) + (\mathrm{A}: (\mathrm{A}: \mathfrak {a})) = \mathrm{A}.
\]

\exercisesection{2}

\begin{enumerate}
\def\labelenumi{\arabic{enumi}.}
\tightlist
\item
  \begin{enumerate}
  \def\labelenumii{(\alph{enumii})}
  \tightlist
  \item
    Show that, if \(\mathfrak{a}\) is an ideal \(\neq 0\) in a Dedekind domain \(\mathbf{A}\) , the ring \(\mathbf{A} / \mathfrak{a}\) is a principal ideal ring (Algebra, Chapter VII, § 1, Exercise 5) (note that \(\mathbf{A} / \mathfrak{a}\) is semi-local and argue as in Proposition 1 of no. 2). Deduce again that every fractional ideal of \(\mathbf{A}\) is generated by at most two elements (cf.~Exercise 11 (a)).
  \end{enumerate}
\end{enumerate}

\begin{enumerate}
\def\labelenumi{(\alph{enumi})}
\setcounter{enumi}{1}
\tightlist
\item
  In the polynomial ring \(\mathbf{A} = k[\mathbf{X},\mathbf{Y}]\) over a field \(k\) , let \(\mathfrak{m}\) be the ideal \(\mathbf{AX} + \mathbf{AY}\) . Show that for every integer \(n\) the minimum number of generators of the ideal \(\mathfrak{m}^n\) is \(n + 1\) .
\end{enumerate}

\begin{enumerate}
\def\labelenumi{\arabic{enumi}.}
\setcounter{enumi}{1}
\tightlist
\item
  \begin{enumerate}
  \def\labelenumii{(\alph{enumii})}
  \tightlist
  \item
    Let \(\mathfrak{a}\) , \(\mathfrak{b}\) be two integral ideals of a Dedekind domain \(\mathbf{A}\) ; show that there exists an integral ideal \(\mathfrak{c}\) which is prime to \(\mathfrak{b}\) and such that \(\mathfrak{ac}\) is principal (use Proposition 9 of §1, no. 5).
  \end{enumerate}
\end{enumerate}

\begin{enumerate}
\def\labelenumi{(\alph{enumi})}
\setcounter{enumi}{1}
\tightlist
\item
  Let \(\mathfrak{a}, \mathfrak{b}\) be two integral ideals of \(\mathbf{A}\) ; show that there exists \(x \neq 0\) in the field of fractions of \(\mathbf{A}\) such that \(x\mathfrak{a}\) is an integral ideal of \(\mathbf{A}\) which is prime to \(\mathfrak{b}\) (same method). Deduce that the module \(\mathfrak{a}/\mathfrak{ab}\) is isomorphic to \(\mathbf{A}/\mathfrak{b}\) .
\end{enumerate}

\begin{enumerate}
\def\labelenumi{\arabic{enumi}.}
\setcounter{enumi}{2}
\item
  Let A be a Dedekind domain, K its field of fractions, P the set of prime ideals \(\neq 0\) of A and for all \(p \in P\) let \(v_{p}\) be the corresponding essential valuation on K. For every sub-A-module M of K and all \(p \in P\) we write \(v_{p}(M) = \inf v_{p}(x)\) (taken in \(\overline{\mathbf{R}}\) ); if \(M \neq 0\) , \(v_{p}(M) < +\infty\) for all \(p \in P\) and \(v_{p}(M) \leqslant 0\) for almost all \(p \in P\) . If M, N are two sub-A-modules of K, show that the relation \(M \subset N\) is equivalent to \(v_{p}(N) \leqslant v_{p}(M)\) for all \(p \in P\) (use Proposition 9 of §1, no. 5). Conversely, for every family \((v_{p})_{p \in P}\) of elements equal to an integer or \(-\infty\) and such that \(v_{p} \leqslant 0\) for almost all \(p \in P\) , there exists a unique sub-A-module M of K such that \(v_{p}(M) = v_{p}\) for all \(p \in P\) .
\item
  Let A be a Dedekind domain and K its field of fractions. If L is a subfield of K such that A is integral over \(A \cap L\) , then \(A \cap L\) is a Dedekind domain (show first that \(A \cap L\) is a Krull domain and then that every prime ideal of \(A \cap L\) is maximal, using Chapter V, §2, no. 1, Theorem 1).
\item
  \begin{enumerate}
  \def\labelenumii{(\alph{enumii})}
  \tightlist
  \item
    Let k be a field, K the field \(k(\mathbf{X}, \mathbf{Y})\) of rational functions in two indeterminates over k and A the polynomial ring K{[}Z{]} in one indeterminate, which is a principal ideal domain. Let L be the subfield \(k(\mathbf{Z}, \mathbf{X} + \mathbf{Y}\mathbf{Z})\) of \(k(\mathbf{X}, \mathbf{Y}, \mathbf{Z})\) ; show that \(A \cap L\) is not a Dedekind domain (prove that there are prime ideals which are not maximal and \(\neq 0\) in this ring).
  \end{enumerate}
\end{enumerate}

\begin{enumerate}
\def\labelenumi{(\alph{enumi})}
\setcounter{enumi}{1}
\tightlist
\item
  Let A be a Dedekind domain which is not a principal ideal domain and for which the ideal class group C(A) is finite * (the ring of integers of a finite algebraic extension of Q has the latter property). * Let \((\mathfrak{a}_{j})_{1 \leqslant j \leqslant r}\) be a representative system of the group I(A) of fractional ideals \(\neq 0\) modulo the subgroup of fractional principal ideals, consisting of integral ideals of A and let \(p_{i} (1 \leqslant i \leqslant s)\) be the prime ideals dividing at least one of the \(a_{j}\) . Let S (resp. T)
\end{enumerate}

denote the multiplicative set of elements \(x \in \mathbf{A}\) such that \(v_{\mathfrak{p}_i}(x) = 0\) for all \(i\) (resp. consisting of 1 and the \(x \in \mathbf{A}\) such that \(v_{\mathfrak{q}}(x) = 0\) for all the prime ideals \(\mathfrak{q}\) distinct from the \(\mathfrak{p}_i\) and \(v_{\mathfrak{p}_i}(x) \geqslant 1\) for all \(i\) ; the hypothesis implies that \(\mathbf{T}\) is not reduced to 1). Show that \(S^{-1}\mathbf{A}\) and \(T^{-1}\mathbf{A}\) are principal ideal domains and that \(A = (S^{-1}A) \cap (T^{-1}A)\) .

\begin{enumerate}
\def\labelenumi{\arabic{enumi}.}
\setcounter{enumi}{5}
\item
  Show that in a Dedekind domain the notions of primary ideal, irreducible ideal, primal ideal (Chapter IV, § 2, Exercise 33), quasi-prime ideal (Chapter IV, § 2, Exercise 34) and power of a prime ideal are identical.
\item
  Let A be a Noetherian domain. Show that the following properties are equivalent:
\end{enumerate}

\((\alpha)\) A is a Dedekind domain.

(β) For every maximal ideal m of A, there exists no ideal a distinct from m and \(m^{2}\) such that \(m^{2} \subset a \subset m\) .

\((\gamma)\) For every maximal ideal \(\mathfrak{m}\) of \(\mathbf{A}\) , the set of primary ideals for \(\mathfrak{m}\) is totally ordered by inclusion.

(δ) For every maximal ideal m of A, every primary ideal for m is a product of prime ideals.

(Prove that each of the properties \((\gamma)\) and \((\delta)\) implies \((\beta)\) ; then show that \((\beta)\) implies that \(A_{m}\) is a field or a discrete valuation ring (cf.~Chapter VI, § 3, no. 5, Proposition 9).)

Give an example of a local integral domain satisfying properties \((\beta)\) , \((\gamma)\) and \((\delta)\) , which is not a valuation ring (cf.~Chapter VI, § 3, Exercise 7).

¶ 8. Let A be an integral domain in which every prime ideal ≠0 is invertible. Show that A is a Dedekind domain. (Prove first that every prime ideal p' ≠ 0 of A is maximal, noting that, if p is a prime ideal such that p ≠ p' and p ⊃ p', then p': p = p' (Chapter II, § 1, Exercises 8 (b)) and on the other hand that p': p = p'p⁻¹, whence a contradiction is deduced. Deduce that A is Noetherian (Chapter II, § 1, Exercise 6 (b)) and finally apply Chapter II, § 5, no. 6, Theorem 4 to show that A\_m is a field or a discrete valuation ring for every maximal ideal m of A.)

\begin{enumerate}
\def\labelenumi{\arabic{enumi}.}
\setcounter{enumi}{8}
\tightlist
\item
  Let A be an integral domain.
\end{enumerate}

\begin{enumerate}
\def\labelenumi{(\alph{enumi})}
\item
  Let \((\mathfrak{p}_i)_{1\leqslant i\leqslant m},(\mathfrak{p}_j')_{1\leqslant j\leqslant n}\) be two finite families of invertible prime ideals such that \(\mathfrak{p}_1\mathfrak{p}_2\dots \mathfrak{p}_m = \mathfrak{p}_1'\mathfrak{p}_2'\dots \mathfrak{p}_n'\) . Show that \(m = n\) and that there exists a permutation \(\pi\) of \([1,n]\) such that \(\mathfrak{p}_i' = \mathfrak{p}_{\pi (i)}\) for all \(i\) such that \(1\leqslant i\leqslant n\) .
\item
  Let \(\mathfrak{p}\) be a prime ideal of \(\mathbf{A}\) and \(a\) an element of \(\mathbf{A} - \mathfrak{p}\) ; if \(\mathfrak{p} \subset \mathbf{A}a + \mathfrak{p}^2\) , show that \(\mathfrak{p} = \mathfrak{p}(\mathbf{A}a + \mathfrak{p})\) ; deduce that, if \(\mathfrak{p}\) is invertible, then \(\mathbf{A}a + \mathfrak{p} = \mathbf{A}\) .
\item
  Show that, if every ideal of A is a product (not necessarily a priori unique) of prime ideals of A, A is a Dedekind domain. (Show first that every invertible prime ideal p is maximal, considering an element \(a \in A - p\) and decompositions of \(Aa + p\) and \(Aa^{2} + p\) as products of prime ideals and applying (a)
\end{enumerate}

in the ring A/p to show that \(Aa^{2} + p = (Aa + p)^{2}\) ; then use (b). Prove then that every prime ideal \(p \neq 0\) is invertible, by decomposing as a product of prime ideals a principal ideal Ab where \(b \in p\) , observing that the factors of this product are invertible and, by applying the above, show that p is necessarily equal to one of these factors. Conclude with the aid of Exercise 8.)

¶ 10. Let A be a ring in which every ideal is a product of prime ideals.

\begin{enumerate}
\def\labelenumi{(\alph{enumi})}
\item
  Show that, for every prime ideal p of A, A/p is a Dedekind domain (Exercise 9).
\item
  Show that the set of minimal prime ideals \(\mathfrak{p}_i\) of \(\mathbf{A}\) is finite (decompose (0) as a product of prime ideals).
\item
  Suppose that \(\mathbf{A} / \mathfrak{p}_i\) is not a field. Show that, if \(y \in \mathfrak{p}_i\) , then, for all \(x \in \mathbf{A} - \mathfrak{p}_i\) , there exists \(z \in \mathfrak{p}_i\) such that \(y = zx\) . (Consider in A the decompositions as products of prime ideals of Ax and Ax + Ay and the corresponding decompositions in \(\mathbf{A} / \mathfrak{p}_i\) .) Deduce that, for all non-zero \(y \in \mathfrak{p}_i\) , \(Ay = \mathfrak{p}_i\) (consider \(\mathfrak{p}_i / Ay\) , decomposing Ay as a product of prime ideals); show finally that \(\mathfrak{p}_i = \mathfrak{p}_i^2\) and therefore that there exists in \(\mathfrak{p}_i\) an idempotent \(e_i\) such that \(\mathfrak{p}_i = Ae_i\) (cf.~Chapter II, § 4, Exercise 15). Then A is the direct composition of the ring \(Ae_i\) and the ring \(\mathbf{A}(1 - e_i)\) , isomorphic to the Dedekind domain \(\mathbf{A} / \mathfrak{p}_i\) .
\item
  Suppose now that all the \(p_{i}\) are maximal, so that A is a direct composition of rings of the form \(A/p_{i}^{v_{i}}\) (Chapter II, § 1, no. 2, Proposition 5) and we may therefore confine our attention to the case where A is primary or where the unique prime ideal p of A is nilpotent. Show that in this case the hypothesis implies that A is a principal ideal ring (Algebra, Chapter VII, § 1, Exercises 5 and 6).
\item
  Conclude from (c) and (d) that A is the product of a finite number of Dedekind domains and a principal ideal ring.
\end{enumerate}

\begin{enumerate}
\def\labelenumi{\arabic{enumi}.}
\setcounter{enumi}{10}
\tightlist
\item
  Let K be a field and \((v_{\iota})_{\iota\in I}\) a family of discrete valuations on K satisfying conditions \((\mathbf{AK}_{\mathbf{I}})\) and \((\mathbf{AK}_{\mathbf{III}})\) of §1, no. 3 and such that, for every integer r, every family \((n_{h})_{1\leqslant h\leqslant r}\) of integers \(\geqslant0\) , every family \((\iota_{h})_{1\leqslant h\leqslant r}\) of distinct elements of I and every family \((a_{h})_{1\leqslant h\leqslant r}\) of elements of K, there exists \(x\in K\) such that \(v_{\iota_{h}}(x-a_{h})\geqslant n_{h}\) for \(1\leqslant h\leqslant r\) and \(v_{\iota}(x)=0\) for ι distinct from the \(\iota_{h}\) . Show that the intersection A of the valuation rings of the \(v_{\iota}\) is a Dedekind domain whose field of fractions is K and that \((v_{\iota})_{\iota\in I}\) is the family of essential valuations of A. (Use Exercise 7 of §1; then prove that two distinct prime ideals of height 1 are relatively prime and deduce that every prime ideal \(\neq0\) of A is maximal.)
\end{enumerate}

¶ 12. Let A be an integral domain and K its field of fractions. Show that the following properties are equivalent:

\((\alpha)\) For every prime ideal \(\mathfrak{p}\) of A, the local ring \(A_{\mathfrak{p}}\) is a valuation ring.

(β) For every maximal ideal m of A, the local ring \(A_{m}\) is a valuation ring.

(γ) Every finitely generated ideal \(\neq0\) in A is invertible (and therefore the finitely generated fractional ideals \(\neq0\) form a group).

(δ) Every torsion-free finitely generated A-module is projective.

(ε) For all \(x \in \mathbf{K}\) , the fractional ideal \(\mathbf{A} + \mathbf{A}x\) is invertible.

(ζ) For all \(x \neq 0\) in \(\mathbf{K}\) , there exists \(y \in \mathbf{K}\) such that \(y \equiv 0 \pmod{1}\) , \(y \equiv 0 \pmod{x}\) and \(y \equiv 1 \pmod{1 - x}\) .

\((\eta)\) If \((a_{i})\) is a finite family of ideals of A, for the system of congruences \(x \equiv c_{i} \pmod{a_{i}}\) to have a solution it is necessary and sufficient that \(c_{i} \equiv c_{j} \pmod{(a_{i} + a_{j})}\) for every pair of indices i, j (``Chinese remainder theorem'').

(θ) If \(a, b, c\) are three ideals of \(A\) , then \(a \cap (b + c) = a \cap b + a \cap c\) .

\begin{enumerate}
\def\labelenumi{(\roman{enumi})}
\item[(ι)]
  If \(a, b, c\) are three ideals of \(A\) , then \(a + (b \cap c) = (a + b) \cap (a + c)\) .
\item[(κ)]
  A is integrally closed and every finitely generated ideal of A is divisorial.
\end{enumerate}

(λ) A is integrally closed and for \(z \neq 0\) in \(\mathbf{K}\) there exist \(x, y\) in \(\mathbf{A}\) such that \(z = x + yz^2\) .

(μ) A is integrally closed and, if a, b, c are three finitely generated fractional ideals \(\neq0\) , the relation ab = ac implies b = c.

\begin{enumerate}
\def\labelenumi{(\alph{enumi})}
\setcounter{enumi}{21}
\tightlist
\item
  For every finitely generated A-module M, the torsion submodule of M is a direct factor of M.
\end{enumerate}

(σ) Every ring B such that A ⊂ B ⊂ K is an integrally closed domain.

Then A is called a Prüferian domain (or a Prüfer domain).

(To prove the equivalence of \((\alpha)\) , \((\beta)\) , \((\gamma)\) and \((\delta)\) , use Chapter II, § 5, no. 2, Theorem 1 and no. 6, Theorem 4. To prove that \((\varepsilon)\) implies \((\gamma)\) , use the identity \((a + b)(b + c)(c + a) = (a + b + c)(ab + bc + ca)\) between three fractional ideals in an integral domain. To prove that \((\zeta)\) implies \((\eta)\) , argue by induction on the number of \(a_i\) . The equivalence of \((\eta)\) , \((\theta)\) and \((\iota)\) has been shown in Algebra, Chapter VI, § 1, Exercise 25. Prove that \((\lambda)\) implies \((\varepsilon)\) by noting that \((\lambda)\) implies the relation

\[
x (\mathbf {A} + \mathbf {A} z) \subset z (\mathbf {A} + \mathbf {A} z),
\]

using Exercise 6 (b) of § 1 and noting that

\[
(\mathbf {A} + \mathbf {A} z) (\mathbf {A} y + \mathbf {A} x / z) = \mathbf {A}.
\]

To prove that \((\mu)\) implies \((\lambda)\) noting that always

\[
z (\mathrm{A} + \mathrm{Az}) \subset (\mathrm{A} + \mathrm{Az} ^ {2}) (\mathrm{A} + \mathrm{Az}).
\]

Prove that \((\kappa)\) implies \((\lambda)\) noting that every fractional principal ideal \(A t\) which contained \(A\) and \(A z^2\) also contains \(A z\) . To prove that \((\nu)\) implies \((\gamma)\) , consider, for a finitely generated integral ideal \(\mathfrak{a} \neq 0\) of \(A\) , a non-zero element \(c \in \mathfrak{a}(A: \mathfrak{a})\) and apply Exercise 32 of §1 to the ideal \(\mathfrak{b} = c \mathfrak{a}\) . To prove that \((\sigma)\) implies \((\beta)\) , reduce it to the case where \(A\) is a local ring, consider an element \(z \in K - A\) and show that \(z^{-1} \in A\) , noting that the hypothesis implies that \(z \in A[z^2]\) .

¶ 13. (a) In a Prüferian domain A, let \(a\) , \(b\) be two finitely generated ideals; show that, if \(b \notin a\) , there exists a finitely generated ideal \(c\) such that \(a \subset c\) , \(a \neq c\) and \(bc \subset a\) (consider the ideal \(a + b\) ).

\begin{enumerate}
\def\labelenumi{(\alph{enumi})}
\setcounter{enumi}{1}
\item
  Deduce from (a) that, if a is a finitely generated ideal in a Prüferian domain A, then in the ring A/a every element which is not a divisor of zero is invertible. In particular, a prime ideal of A cannot be finitely generated if it is maximal.
\item
  Show that in a Prüferian domain A two prime ideals p, p' neither of which is contained in the other are relatively prime (consider, for every maximal ideal m of A, the ideal \(pA_{m} + p'A_{m}\) ).
\item
  In a Prüferian domain A let \(\mathfrak{a}\) be a finitely generated ideal. For \(\mathfrak{a}\) to be primal (Chapter IV, § 2, Exercise 33), it is necessary and sufficient that \(\mathfrak{a}\) be contained only in a single maximal ideal of A (use (b)); \(\mathfrak{a}\) is then irreducible and quasi-prime (Chapter IV, § 2, Exercise 34), so that, for finitely generated ideals, these three notions coincide. For \(\mathfrak{a}\) to be primary, it is necessary and sufficient that it be contained only in a single prime ideal (necessarily maximal) \(\mathfrak{m}\) ; for \(\mathfrak{a}\) to be strongly primary (Chapter IV, § 2, Exercise 27), it is necessary and sufficient also that \(\mathfrak{mA}_{\mathfrak{m}}\) be principal.
\end{enumerate}

\begin{enumerate}
\def\labelenumi{\arabic{enumi}.}
\setcounter{enumi}{13}
\item
  Let A be a Prüferian domain. Show that for an A-module M to be flat, it is necessary and sufficient that it be torsion-free (use Exercise 12 (δ)). Deduce that A is a coherent ring (Chapter I, § 2, Exercise 12).
\item
  \begin{enumerate}
  \def\labelenumii{(\alph{enumii})}
  \tightlist
  \item
    If A is a Prüferian domain, A/p is Prüferian for every prime ideal p of A.
  \end{enumerate}
\end{enumerate}

\begin{enumerate}
\def\labelenumi{(\alph{enumi})}
\setcounter{enumi}{1}
\item
  If A is Prüferian, so is \(S^{-1}A\) for every multiplicative subset S of A not containing 0.
\item
  Let K be a field and \((\mathbf{A}_{\lambda})\) a non-empty right directed family of subrings of K. Show that, if the \(A_{\lambda}\) are Prüferian domains, so is their union A (use Exercise 12 (ζ)).
\end{enumerate}

\begin{enumerate}
\def\labelenumi{\arabic{enumi}.}
\setcounter{enumi}{15}
\tightlist
\item
  Let A be a Prüferian domain, K its field of fractions and L a (finite or otherwise) algebraic extension of K; show that the integral closure A' of A in L is a Prüferian domain. (Let \(x \in L\) and let \(a_{0}X^{n} + a_{1}X^{n-1} + \cdots + a_{n}\) be a polynomial in A{[}X{]} of which x is a root; if we write
\end{enumerate}

\[
a _ {0} \mathrm{X} ^ {n} + a _ {1} \mathrm{X} ^ {n - 1} + \dots + a _ {n} = (\mathrm{X} - x) (b _ {0} \mathrm{X} ^ {n - 1} + \dots + b _ {n - 1}),
\]

show that the \(b_{i}\) and the \(b_{i}x\) belong to \(A'\) , using §1, Exercise 12 (a). If \(a = \sum_{i=0}^{n} A a_{i}\) , \(b = \sum_{j=0}^{n-1} B b_{j}\) , show then that the ideal \(b.B a^{-1}\) is the inverse of \(B + B x\) .)

¶ 17. Let A be an integral domain.

\begin{enumerate}
\def\labelenumi{(\alph{enumi})}
\item
  For A to be Bezoutian (§ 1, Exercise 20), it is necessary and sufficient that it be Prüferian and pseudo-Bezoutian (§ 1, Exercise 21).
\item
  For A to be a Dedekind domain, it is necessary and sufficient that it be a Prüferian Krull domain (show that every divisorial ideal of A is finitely generated, using Exercise 12 (κ) and § 1, Exercise 11 (a); deduce that every prime ideal \(\neq 0\) of A is maximal (Exercise 13 (b)), then that every prime ideal of A is finitely generated and conclude with the aid of Chapter II, § 1, Exercise 6).
\item
  For A to be a Dedekind domain, it is necessary and sufficient that it be Prüferian and strongly Laskerian (Chapter IV, § 2, Exercise 28). (Observe that, if A is Prüferian and strongly Laskerian, every maximal ideal m of A is such that \(A_{m}\) is a discrete valuation ring (Chapter VI, § 3, Exercise 8), then that for all non-zero \(x \in A\) there is a product of a finite number of maximal ideals of A contained in the ideal Ax and conclude that A is a Krull domain.)
\end{enumerate}

\begin{enumerate}
\def\labelenumi{\arabic{enumi}.}
\setcounter{enumi}{17}
\tightlist
\item
  Let A be a ring, the union of a right directed family of subrings \(A_{\alpha}\) which are Dedekind domains. Suppose that, for every prime ideal p of \(A_{\alpha}\) , there exists \(\beta \geqslant \alpha\) such that there are at least two distinct prime ideals of \(A_{\beta}\) lying above p.
\end{enumerate}

\begin{enumerate}
\def\labelenumi{(\alph{enumi})}
\item
  Show that the ring of all algebraic integers (the integral closure of \(\mathbf{Z}\) in \(\mathbf{C}\) ) satisfies the above conditions (cf.~Chapter V, § 2, Exercise 6).
\item
  Let \(v\) be a non-improper valuation on \(\mathbf{A}\) , \(z \in \mathbf{A}\) such that \(v(z) > 0\) and let \(\mathfrak{a}\) be the ideal of \(\mathbf{A}\) consisting of the \(x\) satisfying \(v(x) \geqslant v(z)\) . Show that \(\mathbf{A}: \mathfrak{a} = \mathbf{A}\) and hence that \(\mathfrak{a}\) is not invertible, although for every maximal ideal \(m\) of \(\mathbf{A}\) , \(\mathfrak{a}\mathrm{A}_{m}\) is a principal ideal in the valuation ring \(\mathrm{A}_{\mathfrak{m}}\) (cf.~Chapter II, § 5, no. 6, Theorem 4).
\item
  Deduce from (a) and (b) an example of a Prüferian domain which is completely integrally closed but not a Krull domain (cf.~Exercises 15 (c) and 17 (b) and Chapter V, § 1, Exercise 14).
\end{enumerate}

¶ 19. An integral domain A is called pseudo-Prüferian if the set of finitely generated divisors of D(A) (§ 1, Exercise 11) is a group. A pseudo-Prüferian domain is regularly integrally closed (§ 1, Exercise 29). A pseudo-Bezoutian domain (§ 1, Exercise 21) is pseudo-Prüferian; a Prüferian domain is pseudo-Prüferian; a Krull domain is pseudo-Prüferian.

\begin{enumerate}
\def\labelenumi{(\alph{enumi})}
\item
  Show by examples that the converses of the last three assertions are not necessarily true.
\item
  Let \(\theta\) be an irrational number \(>0\) and \(\Gamma\) the group \(\mathbf{R}^2\) ordered by taking the set of ordered pairs \((\alpha, \beta)\) such that \(\alpha \geqslant 0\) and \(\beta \geqslant \theta \alpha\) as the set of positive elements; the ordered group \(\Gamma\) is a complete lattice. Let B be the pseudo-principal domain derived from \(\Gamma\) by the procedure of § 1, Exercise 22 and let A be the subring the intersection of B and the algebra over \(k\) of the subgroup \(\mathbf{Q}^2\) of \(\mathbf{R}^2\) . Show that A is a completely integrally closed domain but is not pseudo-Prüferian (note that, if K is the field of fractions of A and U the group of invertible elements of A, the ordered group \(\mathrm{K}^*/\mathrm{U}\) is isomorphic to the group \(\Gamma \cap \mathbf{Q}^2\) , ordered by the ordering induced by that on \(\Gamma\) ).
\item
  If the domain A is pseudo-Prüferian, show that the polynomial domain A{[}X{]} is pseudo-Prüferian (cf.~§ 1, Exercise 30 (c)).
\item
  Let A be a pseudo-Prüferian domain and K its field of fractions. Show that the integral closure B of A in an algebraic extension L of K is a pseudo-Prüferian domain (method of Exercise 16, using § 1, Exercise 29 (a)).
\end{enumerate}

* (e) Deduce from Exercise 30 (d) of § 1 an example of an increasing sequence \((\mathbf{A}_n)\) of factorial domains with the same field of fractions whose union is not a regularly integrally closed domain (nor a fortiori a pseudo-Prüferian domain) (in the example quoted, take B to be a discrete valuation ring). *

\begin{enumerate}
\def\labelenumi{\arabic{enumi}.}
\setcounter{enumi}{19}
\tightlist
\item
  Let A be a Dedekind domain, K its field of fractions, L a finite algebraic extension of K and B the integral closure of A in L, which is a Dedekind domain. Let f be an ideal of B. For there to exist a ring C such that \(A \subset C \subset B\) and such that f is the conductor of B in A (Chapter V, § 1, no. 5), it is necessary and sufficient that, for every prime ideal \(p'\) of B containing f such that the field \(B/p'\) is isomorphic to \(A/(p' \cap A)\) (prime ideals of residue degree 1), the intersections of f and the transporter f: \(p'\) of \(p'\) in f with A be equal. (Note that, if there exists such a ring C, the conductor of B in the ring \(C_{0} = A + f\) is also equal to f; the existence of C is therefore equivalent to the fact that \(A + f\) contains no ideal of B distinct from f and containing f.~To prove that the latter condition is equivalent to that of the statement, reduce it to the case where A is a discrete valuation ring.)
\end{enumerate}

¶ 21. (a) Let A be an integral domain, f, g two polynomials in A{[}X{]} and h = fg. Let a, b, c denote the ideals of A generated respectively by the coefficients of f, g, h. Show that, if \(\deg(g) = n\) , then \(a^{n+1}b = a^{n}c\) . (Let \(f(X) = \sum_{i=1}^{m} a_{i}X^{i}\) , \(g(X) = \sum_{j=1}^{n} b_{j}X^{j}\) ; for every increasing sequence

\[
\sigma = (i _ {k}) _ {1 \leqslant k \leqslant n + 1}
\]

of integers \(\leqslant m\) and all \(j\) , let

\[
u _ {\sigma , j} = a _ {i _ {1}} a _ {i _ {2}} \dots a _ {i _ {j + 1}} b _ {j} a _ {i _ {j + 2}} \dots a _ {i _ {n + 1}};
\]

take on the set of \(u_{\sigma,j}\) a total ordering such that, for \(j < j'\) , \(u_{\sigma,j} < u_{\tau,j'}\) for all \(\sigma, \tau\) , and, for all \(j\) , \(u_{\sigma,j} \leqslant u_{\tau,j}\) if and only if \(\sigma \leqslant \tau\) in the lexicographical ordering on \([0, m]^{n+1}\) . Then argue by induction on this totally ordered set.)

\begin{enumerate}
\def\labelenumi{(\alph{enumi})}
\setcounter{enumi}{1}
\item
  Deduce from (a) that, if A is Prüferian (Exercise 12), then c = ab.
\item
  Taking A to be the polynomial ring \(\mathbf{Z}[\mathbf{Y}]\) , give an example of two polynomials \(f, g\) of the first degree in \(\mathrm{A}[\mathbf{X}]\) such that \(c \neq ab\) .
\end{enumerate}

¶ 22. (a) Let A be a Noetherian domain each of whose prime ideals ≠0 is maximal, so that every ideal a ≠ 0 may be written uniquely as a product \(\prod_{i} q_{i}\) of primary ideals relative to the distinct prime ideals \(p_{i}\) containing a.

Let \(\mathfrak{a}\) be an invertible ideal in \(\mathbf{A}\) and let \(\mathfrak{b}\) be any ideal \(\neq 0\) in \(\mathbf{A}\) ; show that there exists an ideal \(\mathfrak{c}\) such that \(\mathfrak{b} + \mathfrak{c} = \mathbf{A}\) and \(\mathfrak{ac}\) is principal (Observe that, if \((\mathfrak{p}_i)_{1 \leqslant i \leqslant n}\) is a finite family of distinct maximal ideals of \(\mathbf{A}\) and we write \(\mathfrak{r}_i = \prod_{j \neq i} \mathfrak{p}_j\) , then \(\mathfrak{a} = \sum_i \mathfrak{ar}_i\) and \(\bigcap_i \mathfrak{ar}_i = \mathfrak{ap}_1\mathfrak{p}_2 \ldots \mathfrak{p}_n\) ; using Chapter II, § 1, no. 2, Proposition 6, deduce the existence of an element \(x \in \mathfrak{a}\) such that \(x^{-1}\mathfrak{a} + \mathfrak{b} = \mathbf{A}\) .) Deduce that \(\mathfrak{a}\) is generated by two elements (cf.~Exercise 1).

\begin{enumerate}
\def\labelenumi{(\alph{enumi})}
\setcounter{enumi}{1}
\tightlist
\item
  Let K be a field and A the subring of the ring of formal power series \(K[[T]]\) consisting of the series \(a_{0} + T^{n}P(T)\) , where \(a_{0} \in K\) , \(P(T) \in K[[T]]\) , for a given integer n.~Show that A is a Noetherian local integral domain with a single prime ideal \(m \neq 0\) , but that the smallest cardinal of a system of generators of m is n.
\end{enumerate}

\exercisesection{3}

\begin{enumerate}
\def\labelenumi{\arabic{enumi}.}
\item
  For an integral domain to be factorial, it is necessary and sufficient that there exist a mapping \(x \mapsto s(x)\) of A - \{0\} to N such that \(s(xy) = s(x) + s(y)\) , the relation \(s(x) = 0\) implies that x is invertible in A and finally, for any two elements x, y of A - \{0\}, neither of which divides the other, there exist elements a, b, z, t of A - \{0\} such that \(ax + by = zt\) , \(s(z) < s(x)\) , t being prime to x and to y. (If this condition is fulfilled, show first that every non-zero element of A is a product of extremal elements and then that, for every extremal element p, Ap is a prime ideal; conclude with the aid of Theorem 1 (d)).
\item
  \begin{enumerate}
  \def\labelenumii{(\alph{enumii})}
  \tightlist
  \item
    Let A be an integral domain, the union of a right directed family \((\mathbf{A}_{\lambda})\) of subrings. Suppose that each of the \(A_{\lambda}\) is factorial and that, if \(\lambda \leqslant \mu\) , every extremal element of \(A_{\lambda}\) is extremal in \(A_{\mu}\) . Show that A is a factorial domain whose set of extremal elements is the union of the sets of extremal elements of each of the \(A_{\lambda}\) (cf.~§ 2, Exercise 19 (e)).
  \end{enumerate}
\end{enumerate}

\begin{enumerate}
\def\labelenumi{(\alph{enumi})}
\setcounter{enumi}{1}
\item
  Deduce from (a) that for every factorial domain \(\mathbf{A}\) the polynomial domain \(\mathbf{A}[\mathbf{X}_{\lambda}]_{\lambda \in \mathbf{L}}\) in any family of indeterminates is factorial.
\item
  Let A be a factorial domain such that every domain of formal power series \(A[[X_{1},\ldots,X_{n}]]\) in a finite number of indeterminates is factorial. Show that every domain of formal power series \(A[[X_{\lambda}]]_{\lambda\in L}\) in any family of indeterminates (Algebra, Chapter IV, §5, Exercise 1) is factorial (same method).
\end{enumerate}

\begin{enumerate}
\def\labelenumi{\arabic{enumi}.}
\setcounter{enumi}{2}
\tightlist
\item
  \begin{enumerate}
  \def\labelenumii{(\alph{enumii})}
  \tightlist
  \item
    Let \(A = \bigoplus_{n \geqslant 0} A_n\) be a graded algebra with positive degrees over a field k; suppose that \(A_0 = k\) and that A is a Krull domain. Show that for A to be factorial, it is necessary and sufficient that every graded prime ideal p of A of height l be of the form p = Aa, where a is a homogeneous element (use Exercise 16 of §1). Show that every homogeneous element \(\neq 0\) of A is a product of homogeneous extremal elements.
  \end{enumerate}
\end{enumerate}

\begin{enumerate}
\def\labelenumi{(\alph{enumi})}
\setcounter{enumi}{1}
\tightlist
\item
  Let \(\mathbf{A}\) be a graded \(k\) -algebra satisfying the conditions of (a). Let \(k'\) be an extension of \(k\) and suppose that \(\mathbf{A} \otimes_k k'\) is a factorial domain; show that \(\mathbf{A}\) is factorial. (If a graded ideal \(\mathfrak{a}\) of \(\mathbf{A}\) is such that \(\mathfrak{a} \otimes_k k'\) is principal in \(\mathbf{A} \otimes_k k'\) , show that \(\mathfrak{a}\) is principal; then use (a).)
\end{enumerate}

\begin{enumerate}
\def\labelenumi{\arabic{enumi}.}
\setcounter{enumi}{3}
\tightlist
\item
  \begin{enumerate}
  \def\labelenumii{(\alph{enumii})}
  \tightlist
  \item
    Show that the ring \(\mathbf{A} = \mathbf{Q}[\mathbf{X},\mathbf{Y}] / p\) , where \(p\) is the principal ideal generated by \(\mathbf{X}^2 +\mathbf{Y}^2 -1\) in \(\mathbf{Q}[\mathbf{X},\mathbf{Y}]\) , is a non-factorial Krull domain (if \(x\) is the image of \(\mathbf{X}\) in \(\mathbf{A}\) , show that \(x\) is an extremal element of \(\mathbf{A}\) but that \(\mathbf{A}x\) is not a maximal ideal of \(\mathbf{A}\) ).
  \end{enumerate}
\end{enumerate}

\begin{enumerate}
\def\labelenumi{(\alph{enumi})}
\setcounter{enumi}{1}
\tightlist
\item
  Show that the ring \(\mathbf{A} \otimes_{\mathbf{Q}} \mathbf{Q}(i)\) is factorial (prove that this ring is isomorphic to the quotient of \(\mathbf{Q}(i) [X, Y]\) by the principal ideal ( \(XY - 1\) ); compare with Exercise 3 (b)).
\end{enumerate}

¶ 5. (a) Let k be a Noetherian factorial domain and B the polynomial ring \(k[X_{1},\ldots,X_{n}]\) where \(n\geqslant3\) ; let \(g_{i}\;(0\leqslant i\leqslant r)\) be elements of \(k[X_{3},\ldots,X_{n}]\)

where \(g_0\) is extremal. Writing \(g = \mathbf{X}_1\mathbf{X}_2 - \sum_{i=0} g_i\mathbf{X}_1^i\) , consider the quotient ring \(\mathbf{A} = \mathbf{B}/g\mathbf{B}\) . Show that \(\mathbf{A}\) is factorial. (Let \(\mathbf{S}\) be the multiplicative subset of \(\mathbf{A}\) generated by 1 and the image of \(\mathbf{X}_1\) in \(\mathbf{A}\) ; apply Proposition 3 of no. 4 to \(\mathbf{S}^{-1}\mathbf{A}\) .)

\begin{enumerate}
\def\labelenumi{(\alph{enumi})}
\setcounter{enumi}{1}
\item
  Let \(k\) be a field of characteristic \(\neq 2\) and \(F\) a homogeneous polynomial of the second degree in \(k[\mathbf{X}_1, \ldots, \mathbf{X}_n]\) where \(n \geqslant 5\) , such that the corresponding polynomial function on \(k^n\) is a non-degenerate quadratic form. Show that the domain \(k[\mathbf{X}_1, \ldots, \mathbf{X}_n] / (\mathbf{F})\) is factorial. (Prove first that a homogeneous polynomial \(G\) of the second degree in \(k[\mathbf{X}_1, \ldots, \mathbf{X}_n]\) , such that the corresponding polynomial function is a non-degenerate quadratic form, is extremal for \(n \geqslant 3\) . Then prove the proposition when \(k\) is algebraically closed, using (a); finally pass to the general case with the aid of Exercise 3 (b).)
\item
  If \(F = X_{1}X_{2} - X_{3}X_{4}\) , show that the domain \(k[X_{1}, X_{2}, X_{3}, X_{4}]/(F)\) is not factorial. (Show that the images of the \(X_{i}\) in this ring are extremal elements.)
\end{enumerate}

¶ 6. (a) Let K be an algebraically closed field of characteristic ≠2. Determine the graded prime ideals of the ring

\[
\mathrm{K} [ \mathrm{X}, \mathrm{Y}, \mathrm{Z} ] / (\mathrm{X} ^ {2} + \mathrm{Y} ^ {2} + \mathrm{Z} ^ {2}).
\]

\begin{enumerate}
\def\labelenumi{(\alph{enumi})}
\setcounter{enumi}{1}
\item
  Let \(k\) be an ordered field, \(a, b, c\) elements \(>0\) of \(k\) and \(A\) the ring \(k[X, Y, Z] / (aX^2 + bY^2 + cZ^2)\) . Show that \(A\) is a factorial domain. (Reduce it to the case where \(k\) is a maximal ordered field, using Exercise 3 (b)); then prove, using (a), that every graded prime ideal of \(A\) of height 1 is principal and apply Exercise 3 (a).)
\item
  Let \(k\) be an ordered field, \(a, b\) elements \(> 0\) of \(k\) and \(B\) the ring
\end{enumerate}

\[
k [ \mathrm{X}, \mathrm{Y} ] / (a \mathrm{X} ^ {2} + b \mathrm{Y} ^ {2} + 1).
\]

Show that B is a factorial domain (use (b) and Exercise 17 (a) of § 1.)

\begin{enumerate}
\def\labelenumi{(\alph{enumi})}
\setcounter{enumi}{3}
\tightlist
\item
  Show that the domain \(\mathbf{C} = \mathbf{Q}[\mathrm{X},\mathrm{Y}] / (\mathrm{X}^2 +2\mathrm{Y}^2 +1)\) is factorial, but that the domain \(\mathbf{C}\otimes_{\mathbf{Q}}\mathbf{Q}(i)\) (where \(i^2 = -1\) ) is not factorial.
\end{enumerate}

¶ 7. Let K be a Noetherian factorial domain and F an extremal element of the polynomial ring \(K[X_{1},\ldots,X_{n}]\) ; suppose that, when a weight \(q(i)>0\) ( \(1\leqslant i\leqslant n\) ) is attributed to each \(X_{i}\) , F is isobaric and of weight q\textgreater0. Let A be the domain generated, in an algebraic closure \(\Omega\) of the field of fractions of \(K(X_{1},\ldots,X_{n})\) , by K, the \(X_{i}\) ( \(1\leqslant i\leqslant n\) ) and a root z of the polynomial \(Z^{c}-F\) , where c is an integer prime to q. Show that A is factorial in the two following cases:

\begin{enumerate}
\def\labelenumi{(\arabic{enumi})}
\item
  \(c \equiv 1 \pmod{q}\) ;
\item
  every finitely generated projective K-module is free (which holds for example when K is a field or a principal ideal domain or a local ring). (In the first case, consider the ring of fractions A{[}1/z{]}; show that it is a factorial domain and apply Proposition 3 of no. 4. In the second case, consider an integer d such that \(cd \equiv 1 \pmod{q}\) and let \(z' \in \Omega\) such that \(z = z'^{d}\) ; the domain B = A{[}z'{]} is factorial by virtue of the first case and is a free A-module. Consider B as a graded ring taking \(z'\) to be of weight q, each of the \(X_{i}\) of weight \(cdq(i)\) and reduce it to proving that for two homogeneous elements u, v of A the ideal Au ∩ Av is principal, using Exercise 16 of §1; consider finally the ideal Bu ∩ Bv in B and use Chapter I, §3, no. 6, Proposition 12.) In particular, if K is a field and a, b, c are three integers \textgreater0 which are relatively prime in pairs, the domain
\end{enumerate}

\[
\mathbf {A} = \mathrm{K} [ \mathrm{X}, \mathrm{Y}, \mathrm{Z} ] / (\mathrm{Z} ^ {a} - \mathrm{X} ^ {b} - \mathrm{Y} ^ {c})
\]

is factorial.

¶ 8. Let A be an integral domain and x, y, z three non-zero elements of A, where x is extremal and Ax ∩ Ay = Axy. Let S be the multiplicative set of the \(x^{n}\) ( \(n \geqslant 0\) ) and let B = S \(^{-1}\) A; consider the domains of formal power series A{[}{[}T{]}{]} and B{[}{[}T{]}{]}.

\begin{enumerate}
\def\labelenumi{(\alph{enumi})}
\tightlist
\item
  Let \(i, j, k\) be three integers \(\geqslant 0\) such that \(ijk - ij - jk - ki \geqslant 0\) and \(z^i \in Ax^j + Ay^k\) . Consider in A{[}{[}T{]}{]} the element \(v = xy - z^{i-1}\mathrm{T}\) . Show that there exists an integer \(t > 0\) and a series
\end{enumerate}

\[
v ^ {\prime} = y ^ {t} x ^ {- 1} + b _ {1} x ^ {- 2} \mathrm{T} + \dots + b _ {n - 1} x ^ {- n} \mathrm{T} ^ {n - 1} + \dots
\]

in B{[}{[}T{]}{]} such that \(vv' \in A[[T]]\) (determine the \(b_{n}\) inductively, proceeding in N by intervals of length ij; in the interior of each interval, take

\[
b _ {n + 1} = b _ {n} z ^ {i - 1} x ^ {- 1};
\]

at the end of each interval, use the inequality \(ijk + ij - jk - ki \geqslant 0\) ).

\begin{enumerate}
\def\labelenumi{(\alph{enumi})}
\setcounter{enumi}{1}
\item
  Suppose that \(z^{i-1} \notin Ax + Ay\) . Show that there exists in the ring A{[}{[}T{]}{]} no formal power series of constant term \(y^k\) ( \(k\) an integer \(>0\) ) and which is an element associated with v in B{[}{[}T{]}{]} (calculate the coefficient of T in the product of v and an invertible element of B{[}{[}T{]}{]}).
\item
  Deduce from (a), (b) and Exercise 7 an example of a factorial domain A such that the domain of formal power series A{[}{[}T{]}{]} is not factorial. (In the above notation and assuming that the conditions of (a) and (b) on x, y, z, i, j, k are fulfilled, show that \(vv'\) cannot be a product of extremal factors \(u_{h}\) ( \(1 \leqslant h \leqslant r\) ) in A{[}{[}T{]}{]}; observe that the \(u_{h}\) are formal power series whose constant terms are powers of y. Then consider the ring C = S'\^{}\{-1\}A{[}{[}T{]}{]}, where S' consists of the series whose constant term is in S; B{[}{[}T{]}{]} is the completion of the Zariski ring C; show that \(v' \in C\) and that v and the \(u_{h}\) are extremal in C; obtain finally a contradiction with (b).)
\item
  Deduce from (c) an example of a Noetherian factorial local integral domain A whose completion \(\hat{A}\) is a non-factorial Krull domain.
\end{enumerate}

¶ 9. (a) Let A be a Noetherian domain such that, for every maximal ideal m of A, \(A_{m}\) is a discrete valuation ring. If B is the ring of formal power series \(A[[X_{1},\ldots,X_{n}]]\) , show that, for every maximal ideal n of B, the domain \(B_{n}\) is factorial (consider its completion, using Proposition 8 of no. 9). Deduce that every divisorial ideal of B is a projective B-module.

\begin{enumerate}
\def\labelenumi{(\alph{enumi})}
\setcounter{enumi}{1}
\item
  Let C be a Noetherian ring such that every finitely generated projective C-module is free; show that the ring of formal power series C{[}{[}X{]}{]} has the same property (cf.~Chapter II, § 3, no. 2, Proposition 5).
\item
  Deduce from (a) and (b) that, if A is a principal ideal domain, the domain of formal power series \(A[[X_{1},\ldots,X_{n}]]\) is factorial.
\end{enumerate}

\begin{enumerate}
\def\labelenumi{\arabic{enumi}.}
\setcounter{enumi}{9}
\tightlist
\item
  \begin{enumerate}
  \def\labelenumii{(\alph{enumii})}
  \tightlist
  \item
    A Prüferian (§ 2, Exercise 12) factorial domain is a principal ideal domain.
  \end{enumerate}
\end{enumerate}

\begin{enumerate}
\def\labelenumi{(\alph{enumi})}
\setcounter{enumi}{1}
\tightlist
\item
  A pseudo-Bezoutian (§ 1, Exercise 21) Krull domain is factorial.
\end{enumerate}

\begin{enumerate}
\def\labelenumi{\arabic{enumi}.}
\setcounter{enumi}{10}
\item
  Let K be a field and A the polynomial domain \(K[X,Y]\) , which is factorial; if L is the field \(K(X^{2},Y/X) \subset K(X,Y)\) , show that the domain \(A \cap L\) is not factorial.
\item
  Show Proposition 5 of no. 8 using Chapter III, § 2, no. 8, Corollary 3 to Theorem 1.
\item
  Extend the Corollary to Proposition 7 of no. 8 to the case where the complete Hausdorff local ring A is not an integral domain (use Algebra, Chapter VIII, §6, Exercise 6 (b)).
\item
  Let A be a complete Noetherian local ring whose residue field is of characteristic p \textgreater{} 0. In the ring of formal power series A{[}{[}T{]}{]} consider the elements \(\omega_{n} = (1 - T)^{p^{n}}\) and \(\gamma_{n} = 1 - \omega_{n}\) for every integer n \textgreater{} 0. Show that \(\gamma_{n}\) is, to within a sign, a distinguished polynomial (no. 8); deduce that \(A_{n} = A[[T]]/(\gamma_{n})\) is identified with the algebra over A of the group
\end{enumerate}

\(G_{n} = Z/p^{n}Z\) . Show that the intersection of the principal ideals \((\gamma_{n})\) is reduced to 0; deduce that A{[}{[}T{]}{]} is identified with the inverse limit \(\lim A_{n}\) .

¶ 15. Let K be a field which is complete with respect to the valuation v, A the ring of the valuation, k its residue field and P a polynomial in A{[}X₁, \ldots, Xₙ{]} of total degree d with the following property: there exists an algebraic extension K' of K such that in K'{[}X₁, \ldots, Xₙ{]} P is a product of polynomials of total degree 1. Suppose further that there exist two polynomials Q, R in A{[}X₁, \ldots, Xₙ{]} such that Q is of total degree s and contains a monomial \(aX_1^{s}\) where φ(a) ≠ 0 (φ denoting the canonical homomorphism A → k), R is of degree ≤d - s and \(\overline{\mathbf{P}} = \overline{\mathbf{Q}}\) . \(\overline{\mathbf{R}}\) (notation of Chapter III, § 4). Show that there then exist in A{[}X₁, \ldots, Xₙ{]} two polynomials Q₀, R₀ of respective degrees s, d - s, such that P = Q₀R₀, \(\overline{\mathbf{Q}} = \overline{\mathbf{Q}}_0\) , \(\overline{\mathbf{R}} = \overline{\mathbf{R}}_0\) and Q₀ contains a monomial \(a_0X_1^{s}\) where φ(a) = φ(a₀). (Consider P, Q, R as polynomials in X₁ with coefficients in the ring B, the completion of A{[}X₂, \ldots, Xₙ{]} with respect to the valuation obtained by extending v by the method of Chapter VI, § 10, no. 1, Proposition 2; then apply Hensel's Lemma; finally use the initial hypothesis on P.)

\begin{enumerate}
\def\labelenumi{\arabic{enumi}.}
\setcounter{enumi}{15}
\item
  Let B be a discrete valuation ring whose residue field k is finite and is not a prime field; let \(k_{0}\) be the prime subfield of k and let A be the subring of B consisting of the elements whose classes in the residue field belong to \(k_{0}\) . Let \(\pi\) be a uniformizer of B and \((\theta_{i})_{1\leq i\leq m}\) a system of invertible elements of B such that the classes \(\bar{\theta}_{i}\) mod. \(\pi\) of the \(\theta_{i}\) form a system of representatives of \(k^{*}\) mod. \(k_{0}^{*}\) . Show that the elements \(p_{i}=\theta_{i}\pi\) and the element \(\pi\) are extremal in A and that every element in A is a product of an invertible element and powers of the \(p_{i}\) and \(\pi\) , although A is not integrally closed.
\item
  \begin{enumerate}
  \def\labelenumii{(\alph{enumii})}
  \tightlist
  \item
    Let A be an integral domain; show that in the ring \(A[X_{ij}]\) , where \((\mathbf{X}_{ij})\) is a family of \(n^{2}\) indeterminates \((1 \leqslant i \leqslant n, 1 \leqslant j \leqslant n)\) , the element \(\det(\mathbf{X}_{ij})\) is extremal. (Reduce it to the case where A is a field; observe that the factors of \(\det(\mathbf{X}_{ij})\) would necessarily be homogeneous polynomials and argue by induction on n.)
  \end{enumerate}
\end{enumerate}

\begin{enumerate}
\def\labelenumi{(\alph{enumi})}
\setcounter{enumi}{1}
\tightlist
\item
  Let \(\mathbf{K}\) be an infinite field and \(\mathbf{F}\) a polynomial in \(\mathbf{K}[\mathbf{Y}_1,\dots ,\mathbf{Y}_m]\) , which is also written as \(\mathrm{F(Y)}\) ; for every square matrix \(s = (\alpha_{ij})\) of order \(m\) with elements in K let \(F(s.Y)\) denote the polynomial F, where the element \(\sum_{j=1}^{m}\alpha_{ij}Y_{j}\) has been substituted for each \(Y_{i}\) . Show that, if F is extremal, so is \(F(s.Y)\) for every invertible matrix s. If there exists an integer \(k\geqslant0\) such that
\end{enumerate}

\[
\mathbf {F} (s. \mathbf {Y}) = (\det (s)) ^ {k} \mathbf {F} (\mathbf {Y})
\]

for every invertible matrix \(s\) , \(\mathbf{F}\) is necessarily homogeneous in each of the \(\mathbf{Y}_j\) ;

moreover, if \(\mathbf{F} = \mathbf{GH}\) where \(\mathbf{G}\) and \(\mathbf{H}\) are two polynomials in \(\mathbf{K}[\mathbf{Y}_1,\dots ,\mathbf{Y}_m]\) , there exist two integers \(p\) and \(q\) such that \(p + q = k\) and

\[
\mathbf {G} (s. \mathbf {Y}) = (\det (s)) ^ {p} \mathbf {G} (\mathbf {Y}), \quad \mathbf {H} (s. \mathbf {Y}) = (\det (s)) ^ {q} \mathbf {H} (\mathbf {Y})
\]

for every invertible matrix \(s\) (use (a)).

\begin{enumerate}
\def\labelenumi{(\alph{enumi})}
\setcounter{enumi}{2}
\tightlist
\item
  Let A be a ring which is not reduced to 0; consider the polynomial ring \(\mathbf{A}[\mathbf{X}_{ij}]\) , where the \(\mathbf{X}_{ij}\) are \(n(n + 1)/2\) (resp. \(2n(2n - 1)/2\) ) indeterminates with \(1 \leqslant i \leqslant j \leqslant n\) (resp. \(1 \leqslant i < j \leqslant 2n\) ); let \(U = (\xi_{ij})\) (resp. \(V = (\eta_{ij})\) ) be the square matrix of order \(n\) (resp. \(2n\) ) over \(\mathbf{A}[\mathbf{X}_{ij}]\) such that \(\xi_{ij} = \mathbf{X}_{ij}\) for \(1 \leqslant i \leqslant j \leqslant n\) and \(\xi_{ij} = \mathbf{X}_{ji}\) for \(i > j\) (resp. \(\eta_{ii} = 0\) for \(1 \leqslant i \leqslant 2n\) , \(\eta_{ij} = \mathbf{X}_{ij}\) for \(1 \leqslant i < j \leqslant 2n\) , \(\eta_{ij} = -\mathbf{X}_{ji}\) for \(i > j\) ). Show that \(\det(U)\) (resp. \(\operatorname{Pf}(V)\) ) is an extremal element in \(\mathbf{A}[\mathbf{X}_{ij}]\) (argue as in (b), considering \(\det(s, U, {}^ts)\) and \(\operatorname{Pf}(s, V, {}^ts)\) ).
\end{enumerate}

\begin{enumerate}
\def\labelenumi{\arabic{enumi}.}
\setcounter{enumi}{17}
\tightlist
\item
  Let K be a field and f = g/h an element of the field of rational functions \(\mathrm{K}(\mathrm{U}, \mathrm{V})\) in two indeterminates, where g and h are two relatively prime polynomials of \(\mathrm{K}[\mathrm{U}, \mathrm{V}]\) . Show that in the field of rational functions
\end{enumerate}

\[
\mathrm{K} \left(\mathrm{X} _ {1}, \mathrm{Y} _ {1}, \dots , \mathrm{X} _ {n}, \mathrm{Y} _ {n}\right)
\]

in \(2n\) indeterminates the determinant \(\operatorname{det}(f(\mathbf{X}_i, \mathbf{Y}_j))\) is equal to:

\[
\left(\prod_ {i, j} h \left(\mathrm{X} _ {i}, \mathrm{Y} _ {j}\right)\right) ^ {- 1} \mathrm{V} \left(\mathrm{X} _ {1}, \dots , \mathrm{X} _ {n}\right) \mathrm{V} \left(\mathrm{Y} _ {1}, \dots , \mathrm{Y} _ {n}\right) \mathrm{F} \left(\mathrm{X} _ {1}, \mathrm{Y} _ {1}, \dots , \mathrm{X} _ {n}, \mathrm{Y} _ {n}\right)
\]

where F is a polynomial in \(K[X_{1}, Y_{1}, \ldots, X_{n}, Y_{n}]\) and \(V(X_{1}, \ldots, X_{n})\) is the Vandermonde determinant (Algebra, Chapter III, § 6, no. 4). Consider the particular case where \(f = 1/(U + V)\) (``Cauchy's identity'').

\begin{enumerate}
\def\labelenumi{\arabic{enumi}.}
\setcounter{enumi}{18}
\tightlist
\item
  If \(U\) is a square matrix of order \(n\) , \(\Delta\) its determinant and \(\Delta_p\) the determinant of the \(p\) -th exterior power of \(U\) (Algebra, Chapter III, § 6, no. 3), show that:
\end{enumerate}

\[
\Delta_ {p} = \Delta^ {\binom {n - 1} {p - 1}}
\]

(use Exercise 11 of Algebra, Chapter III, § 6, and Exercise 17 (a) above).

\begin{enumerate}
\def\labelenumi{\arabic{enumi}.}
\setcounter{enumi}{19}
\tightlist
\item
  Let \(\mathbf{A}\) be a factorial domain and \(f = \sum_{k=0}^{n} a_k \mathbf{X}^k\) a polynomial in \(\mathbf{A}[\mathbf{X}]\) : suppose that there exists an extremal element \(p\) in \(\mathbf{A}\) such that:
\end{enumerate}

\begin{enumerate}
\def\labelenumi{(\arabic{enumi})}
\item
  there exists an index \(k \leqslant n\) such that \(a_{k}\) is not divisible by \(p\) but \(a_{i}\) is divisible by \(p\) for \(i < k\) ;
\item
  \(a_0\) is divisible by \(\pmb{p}\) but not by \(p^2\) .
\end{enumerate}

Show that under these conditions one of the irreducible factors of \(f\) in \(\mathbf{A}[\mathbf{X}]\) is of degree \(\geqslant k\) (argue in \((\mathbf{A} / p\mathbf{A})[\mathbf{X}]\) ). Consider the particular case where \(k = n\) (``Eisenstein's irreducibility criterion'').

\begin{enumerate}
\def\labelenumi{\arabic{enumi}.}
\setcounter{enumi}{20}
\tightlist
\item
  Show that in \(\mathbf{Z}[\mathbf{X}]\) the following polynomials are irreducible: \(\mathbf{X}^n - a\) , where one of the prime factors of \(a\) has exponent 1; \(\mathbf{X}^{2^k} + 1\) , (replace \(\mathbf{X}\) by \(\mathbf{X} + 1\) ); \(\mathbf{X}^4 + 3\mathbf{X}^3 + 3\mathbf{X}^2 - 5\) ; \(5\mathbf{X}^4 - 6\mathbf{X}^3 - a\mathbf{X}^2 - 4\mathbf{X} + 2\) . (Use Exercise 20.)
\end{enumerate}

¶ 22. (a) Let k be an ordered field and A the quotient of the polynomial ring k{[}X, Y, Z{]} by the principal ideal \((\mathbf{X}^{2} + \mathbf{Y}^{2} + \mathbf{Z}^{2} - 1)\) ; let x, y, z denote the canonical images of X, Y, Z in A. Show that A is a factorial domain (consider the ring of fractions \(A_{z-1}\) (notation of Chapter II, § 5, no. 1), using no. 4, Proposition 3).

\begin{enumerate}
\def\labelenumi{(\alph{enumi})}
\setcounter{enumi}{1}
\tightlist
\item
  Let \((e_i)_{1 \leq i \leq 3}\) be the canonical basis of \(\mathbf{A}^3\) and let \(\mathbf{M}\) be the quotient of \(\mathbf{A}^3\) by the monogenous sub-A-module \(\mathbf{N}\) generated by \(xe_1 + ye_2 + ze_3\) ; show that \(\mathbf{M}\) is a projective A-module (form a complementary submodule of \(\mathbf{N}\) in \(\mathbf{A}^3\) ). * (c) Show that, if \(k = \mathbf{R}\) , the A-module \(\mathbf{M}\) is not free (identify \(\mathbf{M}\) with a submodule of the module of continuous sections of the fibre bundle of tangent vectors to the unit sphere and use the fact that there exists no continuous field of tangent vectors \(\neq 0\) at every point of the sphere). *
\end{enumerate}

\begin{enumerate}
\def\labelenumi{\arabic{enumi}.}
\setcounter{enumi}{22}
\tightlist
\item
  Let B be a Krull domain, E its field of fractions, \(\Delta\) a derivation of E such that \(\Delta(\mathbf{B}) \subset \mathbf{B}\) , K the subfield of E the kernel of \(\Delta\) and A the Krull domain \(B \cap K\) ; suppose that K is of characteristic p \textgreater{} 0; then \(E^{p} \subset K\) and \(B^{p} \subset A\) , so that B is the integral closure of A in E and the canonical homomorphism \(\vec{i}: C(A) \to C(B)\) is defined ( \(\S 1\) , no. 10). Let U denote the group of invertible elements of B.
\end{enumerate}

\begin{enumerate}
\def\labelenumi{(\alph{enumi})}
\item
  If \(b \in \mathbf{E}\) is such that \(\mathrm{div}_{\mathbf{B}}(b)\) is the canonical image of a divisor in \(\mathbf{D}(\mathbf{A})\) , show that \(\Delta b / b \in \mathbf{B}\) (note that for every prime ideal \(\mathfrak{P}\) of \(\mathbf{B}\) of height 1 there exists \(b' \in \mathbf{K}\) such that \(v_{\mathfrak{P}}(b) = v_{\mathfrak{P}}(b')\) and observe that \(\mathbf{B}_{\mathfrak{P}}\) is invariant under \(\Delta\) ). Let \(\mathbf{L}\) be the additive subgroup of \(\mathbf{B}\) consisting of the \(\Delta b / b\) (``logarithmic derivatives'') which belong to \(\mathbf{B}\) (for \(b \in \mathbf{E}\) or \(b \in \mathbf{B}\) , which amounts to the same, and \(b \neq 0\) ) and let \(\mathbf{L}'\) be the subgroup of \(\mathbf{L}\) consisting of the \(\Delta u / u\) , where \(u \in \mathbf{U}\) . Show that, for every divisor \(d \in \mathbf{D}(\mathbf{A})\) such that the image of \(d\) is a principal divisor in \(\mathbf{D}(\mathbf{B})\) , the class mod. \(\mathbf{L}'\) of \(\Delta b / b\) for all \(b\) such that \(i(d) = \mathrm{div}_{\mathbf{B}}(b)\) depends only on the class of \(d\) in \(\mathbf{C}(\mathbf{A})\) and derive a canonical injective homophism \(\phi\) of \(\operatorname{Ker}(\bar{i})\) to \(\mathbf{L} / \mathbf{L}'\) .
\item
  Show that, if \(\Delta(B)\) is contained in no prime ideal of \(B\) of height 1 and \([E: K] = p\) , \(\phi\) is bijective. (Reduce it to showing that, if \(b \in E\) is such that \(\Delta b / b \in B\) and \(\mathfrak{P}\) is a prime ideal of \(B\) of height 1 such that \(v_{\mathfrak{P}}(b)\) is not a multiple of \(p\) , then \(e(\mathfrak{P} / \mathfrak{p}) = 1\) , where \(\mathfrak{p} = \mathfrak{P} \cap A\) ; for this, deduce from the hypothesis that, if \(t\) is a uniformizer of \(B_{\mathfrak{P}}\) , then \(\Delta t / t \in B_{\mathfrak{P}}\) , so that \(\mathfrak{P}B_{\mathfrak{P}}\) is invariant under \(\Delta\) and that \(\Delta\) therefore defines by taking quotients a derivation \(\overline{\Delta}\) of the residue field \(k = B_{\mathfrak{P}} / \mathfrak{P}B_{\mathfrak{P}}\) ; show that \(\overline{\Delta} \neq 0\) and deduce that \(f(\mathfrak{P} / \mathfrak{p}) = p\) .)
\item
  Let \(k\) be a field of characteristic 2, B the polynomial ring \(k[\mathbf{X}, \mathbf{Y}, \mathbf{Z}]\) and \(\Delta\) the derivation of \(\mathbf{E} = k(\mathbf{X}, \mathbf{Y}, \mathbf{Z})\) such that \(\Delta(\mathbf{X}) = \mathbf{Y}^4\) , \(\Delta(\mathbf{Y}) = \mathbf{X}^2\) , \(\Delta(\mathbf{Z}) = \mathbf{XYZ}\) ; the field \(\mathbf{K}\) the kernel of \(\Delta\) is such that \([\mathbf{E} : \mathbf{K}] = 4\) . Then \(\Delta(\mathbf{Z}) / \mathbf{Z} \in \mathbf{B}\) , but show that \(\mathrm{div}_{\mathbf{B}}(\mathbf{Z})\) is not the canonical image of a divisor in \(\mathrm{D}(\mathbf{A})\) . (Argue by reductio ad absurdum supposing that \(e(\mathfrak{P} / \mathfrak{p}) = 1\) for \(\mathfrak{P} = \mathbf{B}\mathbf{Z}\) , \(\mathfrak{p} = \mathfrak{P} \cap \mathbf{A}\) ; there would then be in \(\mathbf{A}\) a uniformizer of \(\mathrm{B}_{\mathfrak{P}}\) , necessarily of the form \(a(\mathbf{X}, \mathbf{Y}, \mathbf{Z})\mathbf{Z}\) where \(b = a(\mathbf{X}, \mathbf{Y}, 0) \neq 0\) ; deduce that \(\Delta b / b = -\mathbf{X}\mathbf{Y}\) and obtain a contradiction by calculating \(\Delta(-\mathbf{X}\mathbf{Y})\) .)
\end{enumerate}

\begin{enumerate}
\def\labelenumi{\arabic{enumi}.}
\setcounter{enumi}{23}
\tightlist
\item
  \begin{enumerate}
  \def\labelenumii{(\alph{enumii})}
  \tightlist
  \item
    Let E be a field of characteristic 2, \(\Delta\) a derivation of E and K the subfield of E the kernel of \(\Delta\) ; suppose that \([E:K]=2\) . Show that \(\Delta^{2}=a\Delta\) where \(a\in K\) and that for an element \(t\in E\) to be of the form \(\Delta x/x\) , it is necessary and sufficient that \(\Delta t=at+t^{2}\) .
  \end{enumerate}
\end{enumerate}

\begin{enumerate}
\def\labelenumi{(\alph{enumi})}
\setcounter{enumi}{1}
\tightlist
\item
  Let B be a factorial local integral domain of characteristic 2, m its maximal ideal, E its field of fractions and \(\Delta\) a derivation of E; suppose that the subfield K of E, the kernel of \(\Delta\) , satisfies {[}E: K{]} = 2; moreover, suppose that there exist two elements x, y of m such that \(\Delta x\) and \(\Delta y\) generate the ideal q of B generated by \(\Delta(B)\) . Show then that, if \(t = \Delta z / z\) belongs to q, there exists an invertible element u of B such that \(t = \Delta u / u\) (write \(t = r \Delta x + s \Delta y\) , where r, s are in B, and use (a)).
\end{enumerate}

¶ 25. Let k be a field of characteristic 2, B the ring of formal power series \(k[[X,Y]]\) , E its field of fractions and \(\Delta\) the k-derivation of E defined by \(\Delta(X)=Y^{2j}\) , \(\Delta(Y)=X^{2i}\) (i,j integers \(\geqslant0\) ); the subring A of B consisting of the \(x\in B\) such that \(\Delta x=0\) is the ring of formal power series in \(k[[X,Y,Z]]\) where \(X^{2}\) is substituted for X, \(Y^{2}\) for Y and \(X^{2i+1}+Y^{2j+1}\) for Z.

\begin{enumerate}
\def\labelenumi{(\alph{enumi})}
\item
  Show that the group C(A) contains a vector space over k of dimension \(N(i,j)\) equal to the number of ordered pairs of integers \((a,b)\) such that \(0 \leqslant a < i, 0 \leqslant b < j\) and \((2j + 1)a + (2i + 1)b \geqslant 2ij\) . (Use Exercise 23 (b) and Exercise 24 (a), note that the elements of L are the formal power series \(F \in B\) such that \(\Delta F = F^{2}\) ; attributing to X the weight \(2j + 1\) and to Y the weight \(2i + 1\) , decompose F as an infinite sum of isobaric polynomials; if \(L_{q}\) is the subgroup of L consisting of the \(F \in L\) whose isobaric components are of weight \(\geqslant q\) , \(L/L'\) is isomorphic to the direct sum of the groups \(C_{q}/C_{q+1}\) , where \(C_{q} = L_{q}/(L' \cap L_{q})\) ; calculate these groups for \(q \geqslant 4ij\) , using Exercise 24 (b).)
\item
  Show that the ideal \(AX^{2}\) of A is prime; if \(A' = A[X^{-2}]\) , deduce that C(A') and C(A) are isomorphic (no. 4, Proposition 3). Show that A' is a Dedekind domain (consider the ring \(B' = B[X^{-2}]\) , which is integral over A' and a principal ideal domain, and use Chapter V, § 2, no. 4, Theorem 3). Deduce an example of a Dedekind domain whose ideal class group is infinite.
\end{enumerate}

\begin{enumerate}
\def\labelenumi{\arabic{enumi}.}
\setcounter{enumi}{25}
\tightlist
\item
  Let A be a factorial domain and K its field of fractions. For elements \(i\) ( \(1 \leqslant i \leqslant r\) ) of B = A{[}X \(_{1}\) , \ldots, X \(_{n}\) {]} to be such that the ideal \(\sum_{i=1}^{r} \text{Bf}_i\) is equal to B, it is necessary and sufficient that there exist polynomials \(v_{i}\) ( \(1 \leqslant i \leqslant r\) ) in \(\mathbf{K}[\mathbf{X}_1, \ldots, \mathbf{X}_n]\) such that \(\sum_{i=1}^{r} v_i f_i = 1\) and which satisfy further the following condition: if we write \(v_i = w_i / d\) , where \(d \in \mathbf{A}\) , the polynomials \(w_i\) belong to \(\mathbf{A}[\mathbf{X}_1, \ldots, \mathbf{X}_n]\) and the g.c.d. of the set of coefficients of all the \(w_i\) ( \(1 \leqslant i \leqslant r\) ) is equal to 1, then, for every extremal element \(p\) of A dividing \(d\) , the ideal generated by the classes of the \(f_i\) in the ring \((\mathrm{A} / \mathrm{A}p)[\mathrm{X}_1, \ldots, \mathrm{X}_n]\) is the whole of this ring.
\end{enumerate}

¶ 27. (a) Let K be a field and B the ring generated, in an algebraic closure of the field of rational functions K(U, V, X, Y) in 4 indeterminates, by the polynomial ring K{[}U, V, X, Y{]} and a root z of the polynomial

\[
\mathrm{F} = \mathrm{Z} ^ {7} - \mathrm{U} ^ {5} \mathrm{X} ^ {2} - \mathrm{V} ^ {4} \mathrm{Y} ^ {3}.
\]

Show that \(\mathbf{B}\) is a factorial domain (cf.~Exercise 7).

\begin{enumerate}
\def\labelenumi{(\alph{enumi})}
\setcounter{enumi}{1}
\tightlist
\item
  Let p be the (prime) ideal generated by X, Y, U, V and z in A and write \(C = A_{p}[[T]]\) ; show that C is a Noetherian local integral domain, which is not factorial, but whose associated graded domain \(\text{gr}(C)\) is factorial (use Exercise 8).
\end{enumerate}
